% Standalone English manuscript; embedded bibliography; no external figures.
% Variant: 03_Full_Audit_With_Minimax
\documentclass[11pt]{article}

\usepackage[margin=1in]{geometry}
\usepackage{amsmath,amssymb,amsthm,mathtools,bm}
\usepackage{booktabs,tabularx,array,multirow,longtable}
\usepackage{graphicx}
\usepackage{subcaption}
\usepackage{enumitem}
\usepackage{xspace}
\usepackage{placeins}
\usepackage[numbers,sort&compress]{natbib}
\usepackage{xcolor}
\usepackage{microtype}
\usepackage[colorlinks=true,linkcolor=blue,citecolor=blue,urlcolor=blue]{hyperref}
\usepackage[nameinlink,capitalize]{cleveref}

\graphicspath{{figures/}}

\newtheorem{theorem}{Theorem}[section]
\newtheorem{lemma}[theorem]{Lemma}
\newtheorem{proposition}[theorem]{Proposition}
\newtheorem{corollary}[theorem]{Corollary}
\newtheorem{definition}[theorem]{Definition}
\newtheorem{assumption}[theorem]{Assumption}
\newtheorem{condition}[theorem]{Condition}
\newtheorem{example}[theorem]{Example}
\newtheorem{remark}[theorem]{Remark}

\newcommand{\E}{\mathbb E}
\newcommand{\Pp}{\mathbb P}
\newcommand{\R}{\mathbb R}
\newcommand{\N}{\mathbb N}
\newcommand{\Var}{\operatorname{Var}}
\newcommand{\Cov}{\operatorname{Cov}}
\newcommand{\1}{\mathbf 1}
\newcommand{\op}{\mathrm{op}}
\newcommand{\diag}{\operatorname{diag}}
\newcommand{\tr}{\operatorname{tr}}
\newcommand{\argmin}{\operatorname*{arg\,min}}
\newcommand{\argmax}{\operatorname*{arg\,max}}
\newcommand{\dist}{\operatorname{dist}}
\newcommand{\diam}{\operatorname{diam}}
\newcommand{\supp}{\operatorname{supp}}
\newcommand{\sgn}{\operatorname{sgn}}
\DeclarePairedDelimiter{\abs}{\lvert}{\rvert}
\DeclarePairedDelimiter{\norm}{\lVert}{\rVert}
\DeclarePairedDelimiterX{\ip}[2]{\langle}{\rangle}{#1,#2}

\newcommand{\cA}{\mathcal A}
\newcommand{\cB}{\mathcal B}
\newcommand{\cD}{\mathcal D}
\newcommand{\cE}{\mathcal E}
\newcommand{\cF}{\mathcal F}
\newcommand{\cG}{\mathcal G}
\newcommand{\cH}{\mathcal H}
\newcommand{\cL}{\mathcal L}
\newcommand{\cM}{\mathcal M}
\newcommand{\cP}{\mathcal P}
\newcommand{\cS}{\mathcal S}
\newcommand{\cT}{\mathcal T}
\newcommand{\ThetaSet}{\Theta}
\newcommand{\DeltaN}{\Delta_N}
\newcommand{\CPT}{\ensuremath{\mathrm{CPT}}\xspace}
\newcommand{\CPTPG}{\ensuremath{\mathrm{CPT\mbox{-}PG}}\xspace}
\newcommand{\DRCPTPG}{\ensuremath{\mathrm{DRCPT\mbox{-}PG}}\xspace}
\newcommand{\EUPG}{\ensuremath{\mathrm{EU\mbox{-}PG}}\xspace}
\newcommand{\Dir}{\operatorname{Dirichlet}}
\newcommand{\Reg}{\mathcal R}
\newcommand{\Sta}{\mathsf S}
\newcommand{\eps}{\varepsilon}
\newcommand{\etaP}{\eta_{+}}
\newcommand{\etaM}{\eta_{-}}
\newcommand{\dH}{d_{\mathrm H}}
\newcommand{\db}{\mathrm{db}}
\newcommand{\pl}{\mathrm{pl}}
\newcommand{\cf}{\mathrm{cf}}
\newcommand{\as}{\mathrm{as}}


\theoremstyle{definition}
\newtheorem{algorithm}{Algorithm}
\crefname{algorithm}{Algorithm}{Algorithms}
\Crefname{algorithm}{Algorithm}{Algorithms}
\newcommand{\assess}[3]{\par\noindent\textbf{Role.} #1\par\noindent\textbf{Verification.} #2\par\noindent\textbf{Boundary if violated.} #3\par\medskip}
\newcommand{\dependson}[1]{\par\noindent\textbf{Dependencies.} #1\par\smallskip}
\setlength{\emergencystretch}{2em}
\setlist[enumerate]{itemsep=3pt,topsep=5pt}
\title{Dynamic-Reference Cumulative Prospect Policy Gradients:\\Debiased Estimation, Minimax Limits, and Online Portfolio Tracking}
\author{Anonymous Authors}
\date{}
\begin{document}
\maketitle
\begin{abstract}
Cumulative prospect theory (CPT) evaluates uncertain outcomes relative to a reference point. In portfolio learning that reference can evolve along the wealth path, while the market and episode-start state also change. We study a finite-horizon conditional CPT objective with asymmetric reference adaptation, transaction costs, and a factor-parametrized Dirichlet policy. We quantify how perturbations of the initial reference propagate to terminal relative outcomes, CPT values, and policy gradients. A sample-split antithetic estimator shares one outer score trajectory across a full inner sample and its two halves. We prove its weak-bias and level-second-moment rates, then obtain unbiasedness for the frozen simulation target, finite variance, and finite expected trajectory cost by randomizing the level. Simulation error and level truncation are accounted for separately. For the original normalized projected update, we establish an explicit bound on the average squared constrained stationarity residual under objective variation. Window diagnostics and, under a separate residual error bound, distances to moving stationary sets follow as corollaries. Explicit smooth values and compact Dirichlet likelihoods yield reference-sensitivity and objective-smoothness bounds. In a separate exact-trajectory statistical experiment, a realizable portfolio submodel and an expected-budget testing argument establish a matching $\Theta(B^{-1})$ gradient-estimation rate for fixed dimension. This is statistical rate optimality, not optimal portfolio wealth or online regret. Reproducible controlled computations evaluate the specified estimator and update; they do not establish real-market superiority or human behavioral validity.
\end{abstract}

\paragraph{One-sentence contribution.}
We connect the sensitivity of a path-dependent psychological reference, a concrete CPT gradient estimator with unbounded but integrable Dirichlet scores, and constrained online stationarity in a common error budget.

\section{Introduction}\label{sec:introduction}
CPT incorporates reference dependence, asymmetric gain/loss valuation, and nonlinear probability weighting \citep{kahneman1979prospect,tversky1992advances}. These features change the objective being optimized: it is a functional of a trajectory-outcome distribution, rather than an additive expected reward. In portfolio applications, the comparison point may itself change after gains and losses \citep{arkes2008reference,arkes2010cross}. Two wealth paths ending at the same amount can consequently have different terminal relative outcomes.

Existing CPT policy-gradient methods provide the necessary likelihood-ratio viewpoint \citep{prashanth2016cumulative,jie2018stochastic,lepel2024policygradients}. The updated analysis of \citet{lepel2024policygradients} treats finite-sample estimation bias and includes finite-time approximate-stationarity guarantees. We therefore do not position the present paper as the first CPT gradient method or the first analysis of its estimation error. Our question is more specific: how can a continuous portfolio policy be learned when the reference is updated within every candidate trajectory and the conditional objective changes between episodes?

A fixed reference recursion can be incorporated into an augmented Markov state. Its presence alone does not make a market transition kernel nonstationary. We distinguish exogenous changes in the market law from changes in the episode-start context, including wealth, the previous portfolio, and the initial reference. Both affect the conditional target, even when the reference-update rule is fixed. Differentiation of the current target holds this initial context fixed; adaptation through past decisions is not differentiated again.

The statistical difficulty is concrete. A CPT gradient contains a probability-weight derivative evaluated at an unknown survival function. A finite empirical survival substituted into a nonlinear derivative is generally biased. Moreover, a Dirichlet score contains $\log\omega_i$ and is not bounded sample by sample. We specify the outer/inner sample split, prove the antithetic cancellation with a shared outer score, and work with uniform score moments. The resulting randomized estimator is unbiased for the frozen simulator. Accuracy for the true target additionally depends on simulation error.

The optimization difficulty is separate. Projection can eliminate an outward gradient step at the parameter boundary. We analyze the feasible displacement generated by the true normalized gradient, rather than infer stationarity from the unconstrained gradient norm. A one-step inequality connects this residual to objective improvement and direction MSE. Summing it across changing objectives gives the main bound. A model-specific reference-sensitivity result supplies an explicit component of the variation budget.

\paragraph{Contributions.}
\begin{enumerate}
\item We formulate a conditional dynamic-reference CPT portfolio objective and quantify the propagation of an initial-reference perturbation. The explicit smooth values yield a linear sensitivity coefficient, while a separate likelihood calculation derives all required score moments and policy-parameter smoothness.
\item We give the exact sample-split antithetic statistic, establish $O(2^{-\ell})$ weak bias and $O(2^{-2\ell})$ level second moments, and derive unbiased simulation-target estimation, finite expected cost, MSE, and fixed-episode asymptotic normality.
\item We prove an average constrained-stationarity bound for the normalized projected update. The bound separates objective variation from sampling, simulation, and truncation errors. Distance-to-set statements require an additional error-bound condition.
\item We establish a fixed-dimension minimax gradient-estimation rate under an expected raw-trajectory budget. The lower bound uses an actual Bernoulli-return portfolio model, permits adaptive queries and stopping, and matches the cost-calibrated estimator upper bound.
\item We supply executable controlled checks with an analytically tractable CPT gradient and a genuine episodic learning/execution loop. Claims of market outperformance, adoption, or satisfaction are not inferred from these checks.
\end{enumerate}

\paragraph{Organization.}
\Cref{sec:related} positions the work. \Cref{sec:setup} defines the model and assumptions, and \Cref{sec:algorithm} specifies the three algorithms. \Cref{sec:results} states the results. \Cref{sec:experiments} presents controlled evidence and its limits. \Cref{sec:discussion} discusses interpretation. The appendix retains the notation, reference, estimator, online-analysis, rate-analysis, and reproducibility order; every mathematical result is proved there.

\section{Related Work and Positioning}\label{sec:related}
\paragraph{CPT and reinforcement learning.}
CPT-based reinforcement learning and stochastic optimization were developed by \citet{prashanth2016cumulative,jie2018stochastic}. \citet{lepel2024policygradients} derive a first-order CPT gradient, an order-statistic estimator, and asymptotic and finite-time optimization guarantees. Our gradient identity is an adaptation of that foundation, not an independent novelty claim. The additional analysis concerns a specified dynamic-reference objective, continuous Dirichlet likelihood moments, antithetic sampling, and a constrained moving-target residual.

\paragraph{Behavioral alignment beyond expectation.}
CPT gives a mathematical representation of stipulated preferences. Optimizing that representation is not evidence that the stipulated parameters describe a particular person. The parameters in this paper are inputs; their psychological calibration and the acceptability of recommendations are separate empirical questions.

\paragraph{Reference adaptation.}
Evidence and models of asymmetric reference adaptation motivate the recursion \citep{arkes2008reference,arkes2010cross,he2019realization,qin2022realization}. We retain this mechanism and prove a sensitivity estimate that explicitly depends on horizon, adaptation rates, and the value transformation. The estimate controls a component of online objective variation; it does not assert that faster adaptation improves welfare or wealth.

\paragraph{Nonconvex online learning.}
Local nonconvex criteria and projected or proximal methods are established tools \citep{hazan2017efficient,aydore2019dynamic,hallak2021regret}. Our criterion evaluates each episode's constrained residual at its own iterate. It differs from function-value regret and from evaluating all window objectives at a common parameter. The theorem explicitly addresses our criterion and does not claim optimal regret rates.

\paragraph{Debiased and minimax estimation.}
Randomized telescoping follows \citet{rhee2015unbiased}; related unbiased gradients for nested stochastic optimization are studied by \citet{goda2023unbiased}. These principles do not supply the required CPT coupling automatically. We prove the necessary cancellation and moment estimates for the statistic used here. Fixed-episode normal approximation uses standard probability limit theorems \citep{vandervaart1998asymptotic}. The minimax section uses a classical two-point testing method \citep{tsybakov2009introduction}, but verifies its inputs inside a CPT portfolio model and accounts for adaptive stopping under the same expected trajectory budget as the upper bound. Statistical rate optimality does not imply optimal investment performance.

\paragraph{Portfolio learning.}
The action is a long-only allocation to stocks and cash, with factor features and transaction costs \citep{markowitz1952portfolio,fama1993common}. The Dirichlet parametrization enforces the action simplex. Parameter projection is a different operation and remains part of the learning update. Learning concerns the stochastic policy; deterministic mean execution is evaluated separately.

\section{Problem Setup}\label{sec:setup}
\subsection{Nonstationary episodic portfolio MDP}
Episodes are indexed by $k=1,\ldots,K$, have fixed finite horizon $h\ge1$, and begin at date $t_k$. Within an episode, let
\[
 s_u=(m_u,X_u,r_u,\omega_{u-1}),\qquad
 \omega_u\in\Delta_N:=\{\omega\in\R_+^{N+1}:\textstyle\sum_{i=0}^N\omega_i=1\}.
\]
Here $m_u$ contains market factors, $X_u$ is wealth divided by a fixed initial wealth, $r_u$ is the reference in the same units, and asset $0$ is cash. With $R_{0,u+1}=0$ and $c_{\rm tc}\ge0$,
\begin{equation}\label{eq:wealth-update}
 X_{u+1}=X_u\left(1+\omega_u^\top R_{u+1}
 -c_{\rm tc}\norm{\omega_u-\omega_{u-1}}_1\right),\qquad X_1=1.
\end{equation}
The market model includes a conditional return/state kernel, potentially depending on $k$ and within-episode time. Given the state and action it has no direct dependence on the policy parameter. All random variables and conditional kernels are measurable; state and action spaces are standard Borel spaces. Market dynamics and transaction costs must make the declared analysis domain valid.

The information $\cF_k$ contains the observed history, current parameter, initial state, and fitted simulator at the start of episode $k$. Put $z_k=(m_{t_k},X_{t_k},\omega_{t_k-1})$. The true and simulation targets are
\begin{equation}\label{eq:conditional-targets}
 J_k(\theta)=J(P_k,z_k,r_k;\theta),\qquad
 \widetilde J_k(\theta)=J(\widetilde P_k,z_k,r_k;\theta).
\end{equation}
The simulator and initial state are frozen before learning samples are drawn. Every conditional expectation below is a regular conditional expectation given this information. The quantities $\theta_k,M_k$ and the chosen truncation level are $\cF_k$-measurable. A new episode may incorporate newly observed data, but future data are not inserted into the already frozen episode-$k$ simulator.

\subsection{Asymmetric dynamic references}
For fixed $\eta_+,\eta_-\in[0,1]$ define
\begin{equation}\label{eq:reference-update}
 \Phi(r,x)=
 \begin{cases}(1-\eta_+)r+\eta_+x,&x\ge r,\\
 (1-\eta_-)r+\eta_-x,&x<r.
 \end{cases}
 \qquad r_{u+1}=\Phi(r_u,X_{u+1}).
\end{equation}
The common behavioral specification is $\eta_+>\eta_-$. Neither this ordering nor its psychological calibration is a consequence of the optimization theorem. Every simulated trajectory recursively generates its own reference, and
\begin{equation}\label{eq:relative-outcome}
 Y_k(\tau)=X_{k,h}(\tau)-r_{k,h}(\tau),
\end{equation}
where $X_{k,h}=X_{t_k+h}$ and $r_{k,h}=r_{t_k+h}$ denote the trajectory's terminal values, initialized at the saved episode-$k$ context. The initialization $X_1=1$ in \eqref{eq:wealth-update} refers to the beginning of the entire experiment, not a reset at every episode. Copying one realized terminal reference into all simulated trajectories would define a different objective.

\subsection{Regularized cumulative-prospect objective}
Write $y^+=\max(y,0)$ and $y^-=\max(-y,0)$. Fix $\eps_v>0$ and define for $x\ge0$
\begin{equation}\label{eq:value-functions}
 u_{\alpha,\eps_v}(x)=(x^2+\eps_v^2)^{\alpha/2}-\eps_v^\alpha,
 \quad v_+(y)=u_{\alpha_+,\eps_v}(y^+),\quad
 v_-(y)=\lambda u_{\alpha_-,\eps_v}(y^-),
\end{equation}
where $0<\alpha_+,\alpha_-\le1$ and $\lambda>1$. Both magnitudes vanish at zero. They are continuously differentiable as functions of $y$, including zero; no global second derivative in $y$ is needed. Their global Lipschitz constants can be chosen as
\begin{equation}\label{eq:value-lipschitz}
 L_v^+=\alpha_+\eps_v^{\alpha_+-1},\qquad
 L_v^-=\lambda\alpha_-\eps_v^{\alpha_--1}.
\end{equation}
Smoothing is part of the preference specification and affects all positive magnitudes, not only a stated neighborhood. The analysis holds for fixed $\eps_v$; no convergence to the unsmoothed objective is asserted.
For a strictly increasing base weight $w_\pm$ with positive normalization denominator, define
\begin{equation}\label{eq:prob-weights}
 W_\pm(p)=
 \frac{w_\pm(\eps_p+(1-2\eps_p)p)-w_\pm(\eps_p)}
 {w_\pm(1-\eps_p)-w_\pm(\eps_p)},\qquad 0<\eps_p<\tfrac12.
\end{equation}
The running behavioral choice is the Tversky--Kahneman family
\[
 w_\pm(p)=\frac{p^{\beta_\pm}}
 {\{p^{\beta_\pm}+(1-p)^{\beta_\pm}\}^{1/\beta_\pm}},
\]
with parameters selected so that the weight is increasing on the regularized interval. The results apply to any base weight satisfying the stated conditions; the controlled verification uses $w(p)=p^3$. Thus $W_\pm(0)=0$ and $W_\pm(1)=1$. All value and weight parameters are fixed across sample sizes and levels. With different gain/loss weights, normalization can change their relative scaling and is part of the model specification.

Under a true or simulation trajectory law, let
\[
 S_{k,\theta}^{\pm}(z)=\Pp_{k,\theta}\{v_\pm(Y_k)>z\}.
\]
If $0\le v_\pm(Y_k)\le B_v$, the corresponding target is
\begin{equation}\label{eq:cpt-objective}
 J_k(\theta)=\int_0^{B_v}W_+(S_{k,\theta}^+(z))\,dz
 -\int_0^{B_v}W_-(S_{k,\theta}^-(z))\,dz.
\end{equation}
Simulation survivals define $\widetilde J_k$. Strict survival inequalities are used throughout; atoms and tied observations are allowed. In particular $|J_k|,|\widetilde J_k|\le B_J:=B_v$. The payoff functional for a fixed trajectory does not depend explicitly on $\theta$; all parameter dependence is through its probability law.

\subsection{Moving stationary sets and dynamic local regret}
Let $g_k=\nabla J_k(\theta_k)$ and $\widetilde g_k=\nabla\widetilde J_k(\theta_k)$. For fixed $a_0>0$ and $\gamma>0$, set
\begin{equation}\label{eq:normalized-map}
 d(g)=\frac{g}{\max\{\norm{g}_2,a_0\}},\qquad
 Q_k(\theta)=\frac{\Pi_\Theta(\theta+\gamma d(\nabla J_k(\theta)))-\theta}{\gamma}.
\end{equation}
The threshold $a_0$ is distinct from the cash concentration $a_{0,u}$. The constrained stationary set is
\begin{equation}\label{eq:stationary-set}
 \Sta_k=\{\theta\in\Theta:\ip{\nabla J_k(\theta)}{u-\theta}\le0
 \text{ for every }u\in\Theta\}=\{\theta\in\Theta:Q_k(\theta)=0\}.
\end{equation}
The equality follows from the projection variational inequality and positive scaling of the gradient. A continuous objective attains a maximum on compact $\Theta$; its maximizer satisfies this first-order condition, so $\Sta_k$ is nonempty. Continuity of $Q_k$ makes it closed. It may contain saddles and nonglobal stationary points.

Our primary criterion is
\begin{equation}\label{eq:primary-residual}
 \mathcal E_K=\frac1K\sum_{k=1}^K\E\norm{Q_k(\theta_k)}_2^2.
\end{equation}
For fixed $w\ge1$, $0<\rho\le1$, and
$A_k=\sum_{j=0}^{\min(w-1,k-1)}\rho^j$, define the diagnostic
\begin{equation}\label{eq:smoothed-gradient}
 \overline Q_k=A_k^{-1}\sum_{j=0}^{\min(w-1,k-1)}\rho^j
 Q_{k-j}(\theta_{k-j}),
\end{equation}
\begin{equation}\label{eq:dynamic-regret}
 \Reg_{\rho,w}^{Q}(K)=\sum_{k=1}^K\norm{\overline Q_k}_2^2.
\end{equation}
No reverse inference from a small vector average to small individual residuals is used. $Q_k$ is defined using the true gradient, not the realized noisy parameter displacement. Comparisons keep $a_0,\gamma,\Theta$ and the objective specification fixed.

\subsection{Assumptions, diagnostics, and fallback interpretations}\label{subsec:assumptions}
The following are mathematical domain conditions, not properties established by a finite diagnostic sample. Constants are uniform in the episodes, parameters, and true/simulation models being analyzed. Statements as $K\to\infty$ require the same constants for all $K$.

\begin{assumption}[Bounded domain and likelihood regularity]\label{ass:bounded}
The parameter set $\Theta\subset\R^d$ is nonempty, compact, and convex. It is contained in an open parameter neighborhood on which the policy is defined. Features are measurable, have no explicit dependence on $\theta$, and satisfy $\norm{q_{i,u}}_2\le Q_f$. The horizon and asset count are finite and fixed. The initial states and all analyzed true and simulated trajectories satisfy fixed bounds $|X_u|\le B_X$, $|r_u|\le B_r$, and $0\le v_\pm(Y_k)\le B_v$. The feature bound holds on a bounded open neighborhood $O$ of $\Theta$, whose closure is contained in the parameter domain. The initial context is fixed for differentiation, and the only explicit dependence on $\theta$ in the trajectory law is the Dirichlet action density in \eqref{eq:dirichlet-policy}; the conditional market kernels are parameter independent. Lemma~\ref{lem:dirichlet-moments} proves twice dominated likelihood differentiation and uniform fixed-order score moments from these primitive conditions. Denote its score second-moment and expected log-likelihood-Hessian bounds by $M_2$ and $M_H$. These are derived finite constants, not pointwise bounds on the score.
\end{assumption}
\assess{This condition specifies a genuine model domain and justifies likelihood-ratio differentiation and integrability.}{Compact parameter and feature constraints can be checked from the definition and enforced in code. Wealth/value bounds require a model argument, a specified bounded domain, or an explicitly analyzed stopping rule. Score-moment plots diagnose numerical stress but do not prove population bounds.}{Without the required moments or domination, the gradient identity and finite-variance estimator are not asserted. Finite-interval bounds alone do not authorize uniform infinite-time conclusions.}

\begin{assumption}[Fixed regularized CPT]\label{ass:smooth-cpt}
The normalized weights are increasing members of $C^3([0,1])$, with
\[
 C_j:=\max_{\pm}\sup_{p\in[0,1]}|W_\pm^{(j)}(p)|<\infty,
 \qquad j=1,2,3.
\]
All value and weight parameters remain fixed. The values are exactly \eqref{eq:value-functions}, with the fixed positive $\eps_v$ appearing there. Theorem~\ref{thm:gradient-identity} derives a common objective-gradient Lipschitz bound $L=2B_v[C_2M_2+C_1(M_2+M_H)]$ from Lemma~\ref{lem:dirichlet-moments}. Thus objective smoothness is a conclusion of the primitive likelihood conditions in this version. The role of value smoothing is to give linear reference sensitivity, not to authorize a pathwise derivative through portfolio returns.
\end{assumption}
\assess{The derivative bounds justify differentiation of CPT and control the second-order remainder of the multiplier $f(p,I)=W'(p)(I-p)$.}{Endpoint values and derivatives can be checked analytically on the fixed regularized interval. An implementation must use the same weight and derivatives at every level; numerical Hessian checks are diagnostics, not substitutes for a uniform smoothness argument.}{Changing regularization with sample size changes the target and invalidates the stated telescoping identity. Endpoint-singular unregularized weights are outside these bounds.}

\begin{assumption}[Separate objective variation and simulation error]\label{ass:drift}
For $k<K$, define finite measurable quantities
\begin{align}
 v_k&=\sup_{\theta\in\Theta}|J_{k+1}(\theta)-J_k(\theta)|,
 &V_K&=\sum_{k=1}^{K-1}v_k,\label{eq:value-drift}\\
 d_k^J&=\sup_{\theta\in\Theta}
 \norm{\nabla J_{k+1}(\theta)-\nabla J_k(\theta)}_2.&&\label{eq:gradient-drift}
\end{align}
The simulation discrepancy at the queried parameter satisfies
\begin{equation}\label{eq:simulation-bias}
 \norm{\widetilde g_k-g_k}_2\le\beta_k^{\rm sim},
\end{equation}
where $\beta_k^{\rm sim}$ is $\cF_k$-measurable. The expectations used in each bound are finite. No smallness or sublinearity is presumed unless explicitly imposed in an asymptotic statement.
\end{assumption}
\assess{Function-value variation enters telescoping; gradient variation enters set drift; simulation bias enters direction MSE.}{Under a specified market model these quantities can be bounded analytically. With observed data, fitted-model discrepancies and independent gradients provide estimates with their own error; historical observations do not certify future-model accuracy.}{If variation is linear in $K$ or simulation error persists, the theorem permits a nonzero error floor. A within-episode change already represented by $P_k$ is part of the model; an omitted change is simulation error and is not counted a second time.}

\begin{assumption}[Reference-isolated coupling]\label{ass:reference-coupling}
For the reference-sensitivity statements, when $P,z,\theta$ are fixed and only the initial reference is changed, the policy features and market kernel have no direct dependence on that reference or its descendants. Using the same policy and market random numbers therefore produces the same wealth and portfolio paths. The allowed initial-reference interval covers the comparisons used in the variation decomposition.
\end{assumption}
\assess{This permits a pathwise comparison of reference recursions with a shared score and wealth path.}{It is checkable from the state-to-feature map and market kernel. The factors listed in Algorithm~\ref{alg:dirichlet} can exclude $r$ while the learned parameters still respond to it through CPT gradients.}{If features use $X-r$ or directly use $r$, the basic estimation and residual theorems remain available, but the isolated reference-sensitivity and its explicit variation decomposition are not asserted without a new joint coupling argument.}

\begin{assumption}[Residual error bound; optional]\label{ass:error-bound}
For the fixed $a_0,\gamma,\Theta$ in \eqref{eq:normalized-map}, there is a uniform $\kappa<\infty$ such that
\begin{equation}\label{eq:error-bound}
 \dist(\theta,\Sta_k)\le\kappa\norm{Q_k(\theta)}_2,
 \qquad \theta\in\Theta.
\end{equation}
\end{assumption}
\assess{This additional geometric condition turns small residuals into distances to the stationary set.}{It requires a mathematical property of the specified objective. Local multi-start searches can probe the condition but cannot certify it. A neighborhood-only version additionally requires that the iterates and comparison points stay in that neighborhood.}{Without it, the primary residual bound and estimator results remain valid. For example, flat objectives such as $-x^4$ can violate a linear residual error bound near a stationary point.}

\begin{assumption}[Statistical class and exact-trajectory information]\label{ass:oracle}
Fix $d\ge1$, an interior target parameter $\theta_0=0\in\Theta$, finite bounds $\bar h,\bar N\ge1$, and common constants in Assumptions~\ref{ass:bounded} and \ref{ass:smooth-cpt}. Fix the smooth value functions and normalized weights of this paper and require $W_+'(0)>0$. Consider the class $\mathfrak C$ of conditional portfolio models with $1\le h\le\bar h$, $1\le N\le\bar N$ satisfying those common bounds. The horizon, number of assets, features, reference parameters, and starting context are known; the market law is unknown. Constants are chosen to include the one-period model constructed in Appendix~\ref{app:minimax-proof} for a fixed interval of market probabilities. In that submodel the reference rates, return magnitude, and starting wealth are also known.

All information about the unknown market law is obtained from an exact oracle. A call at any chosen $\theta\in\Theta$ returns a fresh complete trajectory under that policy. Calls are conditionally independent given their chosen parameters and model. An estimator may select query parameters adaptively, randomize independently of the model, and stop at a stopping time $N_{\rm call}$ for its observed history. It must satisfy
\[
 \sup_{P\in\mathfrak C}\E_P N_{\rm call}\le B.
\]
There are no free observations, known market-law parameters, or fitted-model side information containing additional observations of the unknown law. The target is $g_P=\nabla J_P(\theta_0)$ in the common Euclidean parameter coordinates. The dimension and all domain constants stay fixed as $B\to\infty$.
\end{assumption}
\assess{This specifies the distributions, parameter, loss, observation experiment, and budget needed for a minimax statement. The nonzero endpoint slope makes the explicit lower-bound submodel locally identifiable through its CPT gradient.}{The oracle contract and counting of every complete trajectory are checkable by design. Exact access to the unknown law is an idealized statistical experiment; a simulator learned from historical data is a different experiment. The positive endpoint slope is checked from the fixed normalized weight.}{With free information about the law, a different query oracle, a deterministic rather than expected budget, growing dimension, or a different class, the stated matching lower and upper bounds do not transfer automatically. If $W_+'(0)=0$, the lower-bound construction below is not claimed to work; the common estimation upper bound still holds.}

There is no assumption postulating the antithetic rates: they are proved for Algorithm~\ref{alg:debiased}. There is also no separate update-alignment assumption. The projection inequality used in the analysis is proved for the actual normalized direction.

\section{Algorithm}\label{sec:algorithm}
\subsection{Dirichlet policy on the portfolio simplex}
For asset $i=0,\ldots,N$, bounded features may include momentum, reversal, volatility, value, a cash indicator, and the previous position. Set
\begin{equation}\label{eq:dir-score}
 z_{i,u}=\theta^\top q_{i,u},\qquad a_{i,u}=e^{z_{i,u}},\qquad A_u=\sum_{i=0}^Na_{i,u},
\end{equation}
\begin{equation}\label{eq:dirichlet-policy}
 \omega_u\sim\pi_\theta(\cdot\mid s_u)=\Dir(a_{0,u},\ldots,a_{N,u}).
\end{equation}
The deterministic recommendation is the policy mean,
\begin{equation}\label{eq:dirichlet-mean}
 \overline\omega_{i,u}=a_{i,u}/A_u.
\end{equation}
Let $\psi_0$ denote the digamma function. The trajectory score is
\begin{equation}\label{eq:explicit-score}
 G_{k,\theta}(\tau)=\sum_{u=t_k}^{t_k+h-1}\sum_{i=0}^N
 a_{i,u}q_{i,u}
 \{\log\omega_{i,u}-\psi_0(a_{i,u})+\psi_0(A_u)\}.
\end{equation}
It is defined almost surely on the open simplex. Its logarithms are unbounded, even though every fixed-order moment required below can be controlled on the declared compact concentration range.

\begin{algorithm}[Dirichlet factor policy for portfolio generation]\label{alg:dirichlet}
\normalfont
\textbf{Input:} state, bounded features, $\theta$, and mode $e\in\{\mathrm{sample},\mathrm{mean}\}$.
\begin{enumerate}
\item Compute $a_{i,u}=\exp(\theta^\top q_{i,u})$ and $A_u=\sum_i a_{i,u}$.
\item If $e=\mathrm{sample}$, draw $\omega_u\sim\Dir(a_{0,u},\ldots,a_{N,u})$; otherwise set $\omega_{i,u}=a_{i,u}/A_u$.
\item Return $\omega_u\in\Delta_N$. Learning samples also record the score contribution in \eqref{eq:explicit-score}.
\end{enumerate}
\end{algorithm}

\subsection{Debiased CPT policy-gradient update}
Learning optimizes $J_k(\pi_\theta)$ or its simulation counterpart. In general this differs from $J_k(\mathrm{mean\mbox{-}policy}_\theta)$. The mean-execution layer supplies recommendations and the next observed context; no transfer of the stochastic-policy guarantee to a deterministic mean-policy objective is assumed.

\begin{algorithm}[Dynamic-reference CPT policy gradient]\label{alg:drcptpg}
\normalfont
\textbf{Input:} $\theta_1\in\Theta$, $X_1=1$, $r_1$, $\omega_0=(1,0,\ldots,0)$, $h$, fixed CPT/reference parameters, $a_0>0$, $\gamma L\le a_0$, fixed $N_0\ge1$ and $1<a<2$, and $\cF_k$-measurable integers $M_k\ge1$ and $\ell_{k,\max}\ge0$; $\ell_{k,\max}=\infty$ is also allowed.
\begin{enumerate}
\item At episode $k$, record $z_k,r_k,\theta_k$ and freeze the fitted simulator using only $\cF_k$.
\item Execute Algorithm~\ref{alg:dirichlet} in mean mode for $h$ steps with $\theta_k$; update actual wealth and reference by \eqref{eq:wealth-update} and \eqref{eq:reference-update}.
\item From the saved episode-start context and simulator, independently repeat Algorithm~\ref{alg:debiased} $M_k$ times. Each repetition redraws its level, outer trajectory, and complete inner sample. This computation can also be performed before or in parallel with step~2.
\item Set $\widehat g_k=M_k^{-1}\sum_{b=1}^{M_k}Z_{k,b}$ and
\begin{equation}\label{eq:parameter-update}
 \theta_{k+1}=\Pi_\Theta\left(\theta_k+
 \gamma\frac{\widehat g_k}{\max\{\norm{\widehat g_k}_2,a_0\}}\right).
\end{equation}
\item Advance to the observed next context and refit only for the next episode. Return parameters, recommendations, wealth/reference paths, and separately labeled diagnostics.
\end{enumerate}
\end{algorithm}
Historical-gradient smoothing is not part of the update. A window in \eqref{eq:dynamic-regret} is a diagnostic window only. Adding stale gradients requires a separate bound on their MSE relative to the current $g_k$.

All the following sampling definitions are conditional on $\cF_k,\theta_k$ and use the frozen simulation law. Suppress $k,b$. An outer trajectory $\tau^{(0)}$ gives $G_0$ and $I_0^\pm(z)=\1\{v_\pm(Y(\tau^{(0)}))>z\}$. Independent inner trajectories $\tau^{(1)},\ldots,\tau^{(n)}$ give
\begin{equation}\label{eq:empirical-survival}
 \widehat S_n^\pm(z)=\frac1n\sum_{i=1}^n
 \1\{v_\pm(Y(\tau^{(i)}))>z\},\qquad
 f_\pm(p,I)=W_\pm'(p)(I-p).
\end{equation}
Define
\begin{equation}\label{eq:T-definition}
 T_n=G_0\left[\int_0^{B_v}f_+(\widehat S_n^+(z),I_0^+(z))\,dz
 -\int_0^{B_v}f_-(\widehat S_n^-(z),I_0^-(z))\,dz\right].
\end{equation}
Write $\mathcal T_\ell=T_{n_\ell}$ for $n_\ell=N_02^\ell$. For $\ell\ge1$, the two half-sample statistics use the first and second halves of this same inner sample and the identical outer trajectory and score:
\begin{equation}\label{eq:antithetic-delta}
 \Delta_0=\mathcal T_0,\qquad
 \Delta_\ell=\mathcal T_\ell-\tfrac12(\mathcal T_{\ell-1}^{(1)}+\mathcal T_{\ell-1}^{(2)}).
\end{equation}
The two half-statistics are not independent unconditionally because they share the outer sample. Only the inner halves are independent, conditional on that outer sample.

\begin{algorithm}[One multilevel debiased CPT gradient replication]\label{alg:debiased}
\normalfont
\textbf{Input:} frozen context/model, $\theta_k$, integer $N_0\ge1$, $1<a<2$, and cap $c\in\N\cup\{\infty\}$.
\begin{enumerate}
\item Put $p_\ell=(1-2^{-a})2^{-a\ell}$. For $c=\infty$ set $q_\ell=p_\ell$; otherwise set $q_\ell=p_\ell/\sum_{j=0}^cp_j$ for $0\le\ell\le c$.
\item Independently draw a level $\mathsf L\sim q$, one outer trajectory, and $N_02^{\mathsf L}$ inner trajectories from the frozen law.
\item Compute $\mathcal T_{\mathsf L}$. If $\mathsf L>0$, compute the two half-sample statistics with the same $G_0,I_0^+,I_0^-$ and form \eqref{eq:antithetic-delta}; if $\mathsf L=0$ use $\Delta_0=\mathcal T_0$.
\item Return $Z^{(c)}=\Delta_{\mathsf L}/q_{\mathsf L}$ and the raw trajectory count $1+N_02^{\mathsf L}$.
\end{enumerate}
\end{algorithm}
All integrals in \eqref{eq:T-definition} are computed by sorting the inner and outer utility values and integrating the resulting step functions exactly. No rank floor depending on $n$, no extra quadrature approximation, and no independently redrawn half-sample score is included. The construction is sample splitting; multi-fold cross-fitting is not presumed.

For comparison, $\widehat g_{n,m}^{\pl}$ uses one inner sample of size $n$ and $m$ independent outer trajectories, independent of that inner sample, and averages \eqref{eq:T-definition}. These $m$ terms share the inner sample. Their variance is analyzed with that dependence retained.

\section{Main Results}\label{sec:results}
\subsection{Dynamic reference creates genuine path dependence}
Define $\eta_{\min}=\min(\eta_+,\eta_-)$ and $\chi_h=(1-\eta_{\min})^h$. The modulus used below is
\begin{equation}\label{eq:value-modulus}
 \Omega(s)=(L_v^++L_v^-)s,\qquad s\ge0.
\end{equation}
It satisfies $|v_+(y)-v_+(y')|+|v_-(y)-v_-(y')|\le\Omega(|y-y'|)$. Write $D_h(s)=\Omega(\chi_h s)$ and $M_1=\sqrt{M_2}$, where $M_2$ is a common upper bound on $\E\norm{G}_2^2$.

% LEAN_TARGET: reference_path_difference_and_sensitivity
\begin{proposition}[Same terminal wealth, different CPT input; reference sensitivity]\label{prop:path-dependence}
Let $L<r_0<C<H$. Consider $A:r_0\to H\to C$ and $B:r_0\to L\to C$, and assume
$C<(1-\eta_+)r_0+\eta_+H$ and $C>(1-\eta_-)r_0+\eta_-L$.
Then
\begin{equation}\label{eq:path-diff}
 r_A-r_B=\eta_+(1-\eta_-)H-\eta_-(1-\eta_+)L+(\eta_--\eta_+)C.
\end{equation}
If this difference is nonzero, the terminal relative outcomes differ. Their CPT values, when each path is a deterministic outcome, also differ under the strictly increasing signed value transformation of this paper.

For any common wealth path and two initial references,
\begin{equation}\label{eq:reference-contraction}
 |r_h-r_h'|\le\chi_h|r_0-r_0'|.
\end{equation}
Under Assumptions~\ref{ass:bounded}, \ref{ass:smooth-cpt}, and \ref{ass:reference-coupling}, for fixed $P,z,\theta$ and admissible $r,r'$,
\begin{align}
 |J(P,z,r;\theta)-J(P,z,r';\theta)|
 &\le C_1D_h(|r-r'|),\label{eq:reference-value-sensitivity}\\
 \norm{\nabla J(P,z,r;\theta)-\nabla J(P,z,r';\theta)}_2
 &\le M_1(C_1+C_2)D_h(|r-r'|).\label{eq:reference-gradient-sensitivity}
\end{align}
The statements also hold for the frozen simulator.
\end{proposition}
\dependson{The path identity and contraction use only \eqref{eq:reference-update}. The value bound uses the common coupling, value modulus, and $C_1$. The gradient bound additionally uses the score identity in \Cref{thm:gradient-identity}; it is proved after establishing that identity, so no circular dependence is used.}
Different arbitrary random outcome distributions need not have different scalar CPT values. Strict contraction requires $\eta_{\min}>0$; when both rates equal one, $Y$ is identically zero, which changes and can flatten the objective rather than improving it.

\subsection{CPT score-gradient identity}
% LEAN_TARGET: cpt_score_gradient_identity
\begin{theorem}[Regularized CPT policy-gradient identity]\label{thm:gradient-identity}
Under Assumptions~\ref{ass:bounded} and \ref{ass:smooth-cpt}, set
\begin{align}\label{eq:psi-def}
 \psi_{k,\theta}(\tau)
 ={}&\int_0^{B_v}W_+'(S_{k,\theta}^+(z))
 \{\1[v_+(Y_k(\tau))>z]-S_{k,\theta}^+(z)\}\,dz\notag\\
 &-\int_0^{B_v}W_-'(S_{k,\theta}^-(z))
 \{\1[v_-(Y_k(\tau))>z]-S_{k,\theta}^-(z)\}\,dz.
\end{align}
Then
\begin{equation}\label{eq:gradient-identity}
 \nabla J_k(\theta)=\E_{k,\theta}[\psi_{k,\theta}(\tau)G_{k,\theta}(\tau)].
\end{equation}
The same identity with simulation survivals and expectation gives $\nabla\widetilde J_k$. In both models,
\begin{equation}\label{eq:gradient-bound}
 \norm{\nabla J_k(\theta)}_2\le B_g:=2B_vC_1\sqrt{M_2}.
\end{equation}
Moreover $J_k$ and $\widetilde J_k$ have $L$-Lipschitz gradients on $\Theta$, where
\begin{equation}\label{eq:derived-smoothness}
 L=2B_v\{C_2M_2+C_1(M_2+M_H)\}.
\end{equation}
Here the likelihood bounds are established in Lemma~\ref{lem:dirichlet-moments}.
\end{theorem}
\dependson{Only the bounded payoff/likelihood conditions and fixed weight regularity are needed. Reference-isolated coupling, drift, and the residual error bound are not needed.}

\paragraph{Why finite empirical survival can be biased.}
The oracle multiplier in \eqref{eq:psi-def} contains the true survival. Replacing it by an empirical survival inside nonlinear $W'$ generally changes its expectation. This observation concerns the explicitly defined statistic \eqref{eq:T-definition}; it is not an assertion that all estimators in the literature have the same bias.

% LEAN_TARGET: antithetic_rates_and_randomized_debiasing
\begin{proposition}[Plug-in bias and debiased estimator]\label{prop:debias}
Under Assumptions~\ref{ass:bounded} and \ref{ass:smooth-cpt} and the sampling construction in Algorithm~\ref{alg:debiased}, let
\[
 H_f=C_3+2C_2,\quad A_b=\frac{B_vH_f\sqrt{M_2}}4,\quad
 A=\frac{A_b}{N_0},\quad C_0=4B_v^2C_1^2M_2,\quad
 D=\frac{3B_v^2H_f^2M_2}{16N_0^2}.
\]
Conditional on $\cF_k$,
\begin{align}
 \norm{\E T_n-\widetilde g_k}_2&\le A_b/n,\label{eq:plugin-bias}\\
 \E\norm{\Delta_0}_2^2&\le C_0,\qquad
 \E\norm{\Delta_\ell}_2^2\le D2^{-2\ell}\quad(\ell\ge1).
 \label{eq:level-second-moment}
\end{align}
Put $r_a=2^{-a}$, $q_a=2^{a-2}$, and
\begin{equation}\label{eq:uniform-db-constant}
 C_{\db}=\frac{C_0+Dq_a/(1-q_a)}{1-r_a},\qquad
 C_{\rm cost}=1+N_0\frac{1-r_a}{1-2r_a}.
\end{equation}
The untruncated output $Z=\Delta_{\mathsf L}/p_{\mathsf L}$, $1<a<2$, satisfies
\begin{equation}\label{eq:db-unbiased}
 \E[Z\mid\cF_k]=\widetilde g_k,
\end{equation}
\begin{equation}\label{eq:db-second-moment}
 \E[\norm{Z}_2^2\mid\cF_k]\le C_{\db},\qquad
 \E[\operatorname{cost}_{\rm traj}(Z)\mid\cF_k]=C_{\rm cost}<\infty.
\end{equation}
With a finite cap $c$ and the probabilities in Algorithm~\ref{alg:debiased},
\begin{equation}\label{eq:cap-bias}
 \E[Z^{(c)}\mid\cF_k]=\E[\mathcal T_c\mid\cF_k],\quad
 \norm{\E Z^{(c)}-\widetilde g_k}_2\le A2^{-c},\quad
 \E\norm{Z^{(c)}}_2^2\le C_{\db}.
\end{equation}
\end{proposition}
\dependson{\Cref{thm:gradient-identity}, $W\in C^3$, the outer score second moment, independent inner samples, and the shared outer trajectory. The two displayed rates are conclusions, not assumptions.}
The trajectory-count cost treats one complete horizon-$h$ path as one call. Arithmetic cost with sorting is also integrable because $\sum_\ell p_\ell n_\ell\log(n_\ell+1)<\infty$. Neither statement places a deterministic bound on one untruncated call.

\subsection{Finite-sample variance and asymptotic distribution}
% LEAN_TARGET: conditional_replication_mse
\begin{theorem}[Finite-sample variance and target MSE]\label{thm:variance}
Under the conditions of \Cref{prop:debias}, for $M$ conditionally independent complete replications at fixed $k,\theta_k$, the untruncated average satisfies
\begin{equation}\label{eq:variance-bound}
 \E[\widehat g_{k,M}\mid\cF_k]=\widetilde g_k,
 \qquad \tr\Var(\widehat g_{k,M}\mid\cF_k)\le C_{\db}/M.
\end{equation}
With cap $c\in\N\cup\{\infty\}$ and $2^{-\infty}=0$, Assumption~\ref{ass:drift} gives
\begin{equation}\label{eq:total-mse}
 \E[\norm{\widehat g_{k,M}^{(c)}-g_k}_2^2\mid\cF_k]
 \le \frac{C_{\db}}M+2(\beta_k^{\rm sim})^2+2A^22^{-2c}.
\end{equation}
The sample-split plug-in average with one shared inner sample and $m$ outer samples satisfies
\begin{equation}\label{eq:plugin-mse}
 \E[\norm{\widehat g_{k,n,m}^{\pl}-g_k}_2^2\mid\cF_k]
 \le\frac{C_n}{n}+\frac{C_0}{m}+\frac{2A_b^2}{n^2}
 +2(\beta_k^{\rm sim})^2,
 \qquad C_n=2B_v^2(C_1+C_2)^2M_2.
\end{equation}
\end{theorem}
\dependson{\Cref{prop:debias} and the stated independence. The true-target terms additionally use \eqref{eq:simulation-bias}; no adjacent-episode drift enters the within-episode variance identity.}

% LEAN_TARGET: fixed_episode_clt
\begin{theorem}[Asymptotic normality of the debiased gradient estimator]\label{thm:clt}
Under the conditions of \Cref{prop:debias}, fix $k,d,\theta_k$, the regularization, and the conditional simulation law. For untruncated iid replications let $\Sigma_k=\Var(Z_k\mid\cF_k)$. For almost every fixed conditional context,
\begin{equation}\label{eq:clt}
 \sqrt M(\widehat g_{k,M}-\widetilde g_k)
 \rightsquigarrow N(0,\Sigma_k).
\end{equation}
If $\Sigma_k$ is positive definite, then
$\sqrt M\Sigma_k^{-1/2}(\widehat g_{k,M}-\widetilde g_k)\rightsquigarrow N(0,I_d)$,
and replacing $\Sigma_k$ by the sample covariance yields the same limit. Finite second moments suffice. Exact simulation permits centering at $g_k$; otherwise it does not.
\end{theorem}
\dependson{\Cref{prop:debias}, iid replication, and for inverse-covariance normalization only, positive definiteness. The result is not a growing-dimension or across-episode limit.}
For a fixed cap, the analogous limit is centered at $\E \mathcal T_c$. Centering a changing-cap estimator at $\widetilde g_k$ requires an additional vanishing $\sqrt M$-scaled bias and a suitable triangular-array argument, neither of which is claimed here.

\subsection{Dynamic local regret}
% LEAN_TARGET: stationary_set_hausdorff_drift
\begin{lemma}[Drift of constrained stationary sets]\label{lem:set-drift}
Under Assumptions~\ref{ass:bounded} and \ref{ass:smooth-cpt}, the gradient-variation part of Assumption~\ref{ass:drift}, and Assumption~\ref{ass:error-bound} for both episodes, with the same $a_0,\gamma,\Theta$,
\begin{equation}\label{eq:set-drift}
 \dH(\Sta_{k+1},\Sta_k)\le\frac{\kappa}{a_0}d_k^J.
\end{equation}
\end{lemma}
\dependson{The residual error bound, gradient variation, and nonexpansiveness of normalization and projection. This lemma is not a premise of the basic online theorem.}

% LEAN_TARGET: average_constrained_stationarity
\begin{theorem}[Dynamic local regret through constrained residuals]\label{thm:dynamic-regret}
Let $\Theta$ be the nonempty compact convex parameter set in Assumption~\ref{ass:bounded}. Suppose $|J_k|\le B_J$, $\norm{\nabla J_k}_2\le B_g$, each $J_k$ is $L$-smooth on $\Theta$, and
\[
 \E[\norm{\widehat g_k-g_k}_2^2\mid\cF_k]\le e_k^2.
\]
For Algorithm~\ref{alg:drcptpg} with fixed $a_0,\gamma>0$, $\gamma L\le a_0$, and $C_g=\max(B_g,a_0)$,
\begin{equation}\label{eq:main-bound}
 \mathcal E_K\le
 \frac{8(2B_J+\E V_K)}{a_0\gamma K}
 +\frac{10C_g}{a_0^3K}\sum_{k=1}^K\E e_k^2.
\end{equation}
Under Assumptions~\ref{ass:bounded}--\ref{ass:drift} and Algorithm~\ref{alg:debiased}, one may take
\begin{equation}\label{eq:algorithm-mse-substitution}
 e_k^2=\frac{C_{\db}}{M_k}+2(\beta_k^{\rm sim})^2
 +2A^22^{-2\ell_{k,\max}}.
\end{equation}
Also $K^{-1}\E\Reg_{\rho,w}^{Q}(K)\le w\mathcal E_K$.

For an explicit reference component, suppose additionally Assumption~\ref{ass:reference-coupling} holds, and define
\begin{equation}\label{eq:nonreference-variation}
 u_k=\sup_{\theta\in\Theta}
 |J(P_{k+1},z_{k+1},r_k;\theta)-J(P_k,z_k,r_k;\theta)|.
\end{equation}
Then the same bound holds with $V_K$ replaced by
\begin{equation}\label{eq:reference-budget}
 U_K+C_1\sum_{k=1}^{K-1}D_h(|r_{k+1}-r_k|),
 \qquad U_K=\sum_{k=1}^{K-1}u_k.
\end{equation}
\end{theorem}
\dependson{\Cref{thm:variance}, the corrected projected-ascent \Cref{lem:projected-ascent}, and value variation. The explicit reference decomposition additionally uses \Cref{prop:path-dependence}. Neither alignment nor an error bound is needed for \eqref{eq:main-bound}.}
The result concerns an average constrained first-order residual. It is not a guarantee of monotone objective or wealth improvement, global CPT maximization, or a fixed recovery time after shocks. All constants may depend on the fixed dimension, horizon, concentration domain, value scale, and regularization. The displayed constants are conservative, not claimed optimal.

% LEAN_TARGET: static_and_sublinear_variation_limit
\begin{corollary}[Static-reference limit]\label{cor:static}
If the complete conditional objective is fixed, $J_k\equiv J$, and
$K^{-1}\sum_{k=1}^K\E e_k^2\to0$, then for fixed $\gamma L\le a_0$,
\begin{equation}\label{eq:static-corollary}
 \frac1K\sum_{k=1}^K\E\norm{Q(\theta_k)}_2^2\longrightarrow0.
\end{equation}
More generally, the same conclusion for $Q_k$ holds when $\E V_K=o(K)$ and average MSE vanishes. Exact untruncated simulation and $M_k\to\infty$ almost surely are sufficient for vanishing average sampling error. A fixed nonzero error bound yields only a possible residual floor.
\end{corollary}
\dependson{\Cref{thm:dynamic-regret} and the stated limiting conditions, with uniform constants.}
To obtain $J_k\equiv J$, fix the market law, episode-start state distribution, reference mechanism, and CPT parameters. A static-reference instance sets $\eta_+=\eta_-=0$ in both real and simulated paths. Observing $r_{k+1}=r_k$ alone is insufficient. No almost-sure convergence to a single parameter is asserted.

% LEAN_TARGET: average_distance_to_stationary_set
\begin{corollary}[Tracking moving first-order stationary sets]\label{cor:set-tracking}
If Assumption~\ref{ass:error-bound} also holds, then
\begin{equation}\label{eq:set-tracking}
 \frac1K\sum_{k=1}^K\E\dist^2(\theta_k,\Sta_k)
 \le\kappa^2\mathcal E_K.
\end{equation}
The right side is bounded by \eqref{eq:main-bound}; under the vanishing conditions of \Cref{cor:static}, it converges to zero.
\end{corollary}
\dependson{Only \Cref{thm:dynamic-regret} and the residual error bound. No inference from a small windowed vector average is used.}

\subsection{Minimax lower bound}
Here ``minimax'' has a statistical meaning: the smallest worst-case mean squared error for estimating one fixed-episode CPT gradient with an expected budget of raw trajectory observations. It does not refer to the maximum attainable CPT value, portfolio wealth, or an online-regret optimum.

% LEAN_TARGET: realizable_cpt_expected_budget_minimax
\begin{theorem}[Fixed-dimension minimax CPT gradient estimation]\label{thm:minimax}
Under Assumption~\ref{ass:oracle}, there exist finite constants $c_{\rm lb},C_{\rm ub}>0$ and $B_0\ge1$, independent of $B$, such that for $B\ge B_0$,
\begin{equation}\label{eq:minimax}
 \frac{c_{\rm lb}}B\le
 \inf_{\widehat g:\,\sup_{P\in\mathfrak C}\E_P N_{\rm call}\le B}
 \sup_{P\in\mathfrak C}\E_P\norm{\widehat g-g_P}_2^2
 \le\frac{C_{\rm ub}}B.
\end{equation}
The infimum is over all measurable procedures with the stated adaptive oracle access. The upper bound is attained by averaging
$M=\lfloor B/C_{\rm cost}\rfloor$ independent untruncated replications of Algorithm~\ref{alg:debiased} at $\theta_0$, with
$B_0=\max(1,2C_{\rm cost})$ and $C_{\rm ub}=2C_{\db}C_{\rm cost}$. Uniform constants are taken over the fixed class. The lower bound holds even when adaptive policy queries and stopping are allowed.
\end{theorem}
\dependson{The upper bound uses \Cref{prop:debias,thm:variance} with exact sampling. The lower bound uses the realizable Bernoulli-return Dirichlet submodel, its nonzero CPT-gradient derivative, the stopped-transcript KL identity, Pinsker's inequality, and a two-point testing reduction, all specified in Appendix~\ref{app:minimax-proof}. Drift, reference-isolated sensitivity, and stationary-set error bounds are not premises.}

The rate is $\Theta(B^{-1})$ for fixed dimension and fixed regularization. Constants can depend on $d$, horizon, features, and regularization. We do not claim a dimension-uniform $d/B$ rate, an optimal constant, or a unique rate-optimal algorithm. In fact, the independently justified plug-in bound with $n$ and $m$ both proportional to $B$ also gives an $O(B^{-1})$ MSE upper bound under exact simulation; the randomized estimator's additional property is exact unbiasedness with finite expected cost. The statistical optimality theorem is not used to prove the online stationarity result.


\section{Experimental Design and Evidence}\label{sec:experiments}
The numerical evidence is a controlled verification of the stated mathematical objects. All reported estimator values are computed from sampled trajectories by Algorithm~\ref{alg:debiased}; all replay updates use Algorithm~\ref{alg:drcptpg}. The evidence does not include a validated A-share dataset or observations of human investors. The final two subsections specify the resulting evidentiary boundary and the measurements needed for those applications. They do not report invented market or user outcomes.

Use $\alpha_+=\alpha_-=0.88$, $\lambda=2.25$, $\eta_+=0.4$, $\eta_-=0.1$, and the normalized base weight $w_\pm(p)=p^3$ with $\eps_p=0.05$. The smooth values in \eqref{eq:value-functions} use $\eps_v=0.01$. The polynomial weight is an analytically tractable member of the declared smooth-weight class, not a fitted behavioral weighting curve. For all randomized estimators $N_0=2$ and $a=1.5$. Random-number seeds and raw episode records are supplied with the source.

\subsection{Behavioral mechanism: path dependence and disposition effect}
Starting from $r_0=X_0=1$, consider wealth paths $(1,1.4,1.2)$ and $(1,0.8,1.2)$. Direct application of \eqref{eq:reference-update} gives terminal references $1.176$ and $1.068$, respectively, hence relative outcomes $0.024$ and $0.132$. These are two feasible no-cost wealth paths with positive gross returns and identical terminal wealth. Their deterministic CPT values differ by strict monotonicity of the signed value function and endpoint normalization. This example evaluates each branch explicitly: the first path's final step is on the gain branch for the chosen parameters, so the special two-branch formula in \eqref{eq:path-diff} is not invoked for these numbers.

This verifies path sensitivity of the mechanism. It does not establish a disposition effect, which is a statement about selling or holding behavior conditional on unrealized gains and losses. Such a claim would require specified liquidation decisions, an appropriate comparison population, and outcome measurements. The current long-only allocation model alone supplies no human-behavior evidence of that kind.

\subsection{Estimator diagnostics: bias, variance, CLT and rate scope}
For an independently tractable one-period instance, take $X_0=r_0=1$, zero transaction costs, a constant positive stock return $r_*=1/3$, and $\theta=0$. Then $U\sim\operatorname{Uniform}(0,1)$, $Y=cU$ with $c=0.2$, and $G=1+\log U$. The true scalar gradient is
\begin{equation}\label{eq:diagnostic-gradient}
 g=\int_0^1c\,v_+'(ct)W_+'(1-t)(-t\log t)\,dt.
\end{equation}
Numerical integration gives $g=0.0643115914$; doubling the integration order from $512$ to $1024$ changes this by $1.804\times10^{-16}$. This is a numerical stability check, not a certified quadrature bound.

For the specified cubic base weight, write $W'(p)=A_0+A_1p+A_2p^2$. The exact finite-inner-sample expectation satisfies
\begin{equation}\label{eq:diagnostic-bias}
 \E T_n-g=B_1/n,\quad
 B_1=A_2\int_0^1 c\,v_+'(ct)t(1-t)(-t\log t)\,dt>0.
\end{equation}
Appendix~\ref{app:experiments} derives this identity. In this instance $B_1=0.0289720186$. The positive bias provides a concrete example for the general bias discussion. At each $n$, Table~\ref{tab:level-diagnostics} uses $1600$ independently sampled outer/inner pairs and their actual coupled half-statistics. Empirical biases can be noisy and may change sign even when the population bias is positive.

\begin{table}[htbp]
\centering\small
\caption{Actual coupled-statistic diagnostics. Parentheses are Monte Carlo standard errors. The analytical bias column is \eqref{eq:diagnostic-bias}; no fitted bias or variance curve is used.}\label{tab:level-diagnostics}
\begin{tabular}{rrrr}
\toprule
$n$ & Analytical bias & Empirical bias (SE) & Mean $\Delta_\ell^2$ (SE)\\
\midrule
4 & $7.243\times10^{-3}$ & $7.087\times10^{-3}$ ($4.033\times10^{-3}$) & $2.309\times10^{-3}$ ($3.260\times10^{-4}$) \\
8 & $3.622\times10^{-3}$ & $3.026\times10^{-4}$ ($3.379\times10^{-3}$) & $4.855\times10^{-4}$ ($6.550\times10^{-5}$) \\
16 & $1.811\times10^{-3}$ & $9.245\times10^{-4}$ ($3.010\times10^{-3}$) & $1.241\times10^{-4}$ ($1.633\times10^{-5}$) \\
32 & $9.054\times10^{-4}$ & $-2.300\times10^{-3}$ ($2.837\times10^{-3}$) & $5.643\times10^{-5}$ ($2.124\times10^{-5}$) \\
64 & $4.527\times10^{-4}$ & $1.621\times10^{-3}$ ($3.062\times10^{-3}$) & $1.344\times10^{-5}$ ($3.352\times10^{-6}$) \\
128 & $2.263\times10^{-4}$ & $-7.934\times10^{-3}$ ($2.445\times10^{-3}$) & $1.462\times10^{-6}$ ($1.823\times10^{-7}$) \\
\bottomrule
\end{tabular}
\end{table}

Table~\ref{tab:batch-diagnostics} uses $8192$ actual randomized replications with cap $c=7$. For each $M$, consecutive nonoverlapping groups of $M$ replications are averaged. The same pool is regrouped across different values of $M$, so rows are not independent experiments. The loss is squared error relative to $g$ in \eqref{eq:diagnostic-gradient}; reported standard errors are the sample standard deviation of group losses divided by the square root of the number of groups. Costs count every outer and inner trajectory. The cap's population target differs from $g$ by $B_1/(2\cdot2^7)=1.132\times10^{-4}$. Consequently these computations do not test the exact unbiasedness of an infinite-level estimator.

\begin{table}[htbp]
\centering\small
\caption{Capped estimator MSE and measured raw trajectory cost per batch.}\label{tab:batch-diagnostics}
\begin{tabular}{rrrrr}
\toprule
$M$ & Groups & Empirical MSE & SE of MSE & Mean trajectory calls\\
\midrule
1 & 8192 & $9.572\times10^{-2}$ & $7.053\times10^{-3}$ & 5.146 \\
4 & 2048 & $2.403\times10^{-2}$ & $1.900\times10^{-3}$ & 20.586 \\
16 & 512 & $6.303\times10^{-3}$ & $6.145\times10^{-4}$ & 82.344 \\
64 & 128 & $1.824\times10^{-3}$ & $2.611\times10^{-4}$ & 329.375 \\
\bottomrule
\end{tabular}
\end{table}

The finite sample tables are compatible with the variance reduction of independent averaging but are not a proof of an asymptotic rate or normal limit. Theorem~\ref{thm:clt} concerns iid untruncated replications at a fixed episode. A finite-cap diagnostic is centered at its own expectation if a normal approximation is used. The minimax statement in Theorem~\ref{thm:minimax} is established analytically for its specified model class and oracle. These finite capped computations neither prove that lower bound nor realize a deterministic-budget version of its untruncated upper-bound algorithm.

\subsection{Synthetic nonstationary market tracking}
We execute a complete $36$-episode replay with $h=1$, one stock and cash, no transaction costs, $\Theta=[-2,2]$, and $X_1=r_1=1$, $\theta_1=0$. Stock returns are $\pm0.04$. The positive-return probability is $0.7$ for episodes $1$--$12$, $0.3$ for episodes $13$--$24$, and $0.6$ for episodes $25$--$36$. At each episode the stochastic-policy learning oracle uses that episode's true distribution, so $\beta_k^{\rm sim}=0$. Mean execution updates actual wealth and reference using a separate return draw. Learning then changes $\theta$ by the specified normalized projected step.

The run uses $M_k=64$, cap $6$, and $a_0=1$. A conservative analytic smoothness bound for this entire finite replay gives $L=177.783288$ and $\gamma=1/(2L)=0.00281241$; Appendix~\ref{app:experiments} derives this choice from a wealth bound and the Beta likelihood. The step is deliberately tied to the theorem's sufficient condition. The resulting small parameter displacement is reported, not presented as evidence of rapid adaptation.

\begin{table}[htbp]
\centering\small
\caption{One executed replay, seed $88177$ for the actual return stream. Residuals and gradients use independent numerical integration of the current true target.}\label{tab:replay}
\begin{tabular}{lr}
\toprule Quantity & Observed value\\\midrule
Mean numerically evaluated $Q_k^2$ & $1.144\times10^{-4}$ \\
Mean squared gradient error & $2.101\times10^{-4}$ \\
Final parameter $\theta_{37}$ & 0.00031455 \\
Final normalized wealth & 1.03335402 \\
Final reference & 1.09220849 \\
Total learning trajectory calls & 11324 \\
\bottomrule
\end{tabular}
\end{table}

The gradient integrator is checked by central differences of the separately integrated objective at $\theta=0.2$ and starting references $0.98,1,1.02$; every difference is below $2\times10^{-7}$. Across the replay, increasing quadrature order from $128$ to $256$ changes a reported objective or gradient by at most $3.902\times10^{-6}$. The residual average in Table~\ref{tab:replay} is a realized numerical diagnostic, not the expectation $\mathcal E_K$ estimated from many independent market replays. The table does not certify the unknown variation budget or a stationary-set error bound. No shock-recovery time, outperformance, or superiority over a baseline is inferred from one path.

\subsection{Semi-synthetic A-share-style replay}
The present evidence supports the estimator and controlled online computation only. No externally validated A-share replay is reported. For a market study to support an additional performance statement, its target must be declared before comparison: the CPT objective of a randomized policy, the realized utility of a deterministic recommendation, and terminal wealth are distinct quantities in this model.

A reproducible market evaluation would record dated prices, the point-in-time investable universe, corporate-action treatment, trading calendars, transaction-cost units, and the history available to each fitted simulator. It would compare dynamic-reference and static-reference CPT variants, an expected-utility variant, and a plug-in version of the same policy under matched raw-trajectory budgets and data splits. Simulators would be fitted only with past data; costs and turnover would be applied to executed mean allocations. Market performance intervals would require repeated independent experimental units or a time-series resampling analysis appropriate to the declared dependence. These requirements describe the evidence needed for that application; no numerical market benefit is included among this paper's conclusions.

\subsection{Heterogeneous retail-agent evaluation}
CPT parameters and reference rates describe a stipulated mathematical preference model. Varying them in a synthetic population can measure sensitivity of recommendations to the stipulation. It cannot establish real-person satisfaction, adoption, trust, or calibrated heterogeneity. Those outcomes require directly observed decisions and a documented measurement protocol.

For a future behavioral evaluation, the relevant estimands would include the proportion of recommendations accepted, signed changes in independently elicited utility, and withdrawal or switching behavior, with prespecified handling of repeated observations from the same person. Parameters estimated from training decisions would be evaluated on held-out decisions. Synthetic-agent scores would be labeled as model-generated outputs and analyzed separately. The theoretical results here neither assume nor establish such empirical behavioral validity.

\section{Discussion}\label{sec:discussion}
The mathematical conclusion is a coupled estimation-and-stationarity statement for a specified conditional CPT objective. The reference recursion makes the distributional payoff path dependent. Under reference-isolated coupling its perturbation contributes an explicit term $C_1D_h(|r_{k+1}-r_k|)$ to objective variation. Independently, the sample-split antithetic construction cancels first-order empirical-survival errors while allowing logarithmically unbounded Dirichlet scores. These two calculations enter the error budget of the normalized projected update. The contribution is this explicit connection and its verified interfaces, rather than the use of CPT as a substitute scalar reward in an otherwise unchanged policy-gradient theorem.

The explicit smooth values yield linear reference sensitivity, and the compact Dirichlet likelihood argument establishes the required moments and parameter smoothness. The statistical optimality statement has a narrower, separate scope: for a fixed-dimensional CPT gradient and the exact fresh-trajectory experiment in Assumption~\ref{ass:oracle}, no procedure improves the worst-case $B^{-1}$ MSE order, and the randomized estimator attains that order. Constants are not claimed uniform in dimension. This does not show uniqueness of the estimator, optimal wealth, or a minimax online-regret rate. Access to a fitted simulator is not free access to the true unknown market law; estimation of that simulator still contributes $\beta_k^{\rm sim}$ in the online theorem.

The online result permits a persistent residual floor if sampling effort is fixed, the simulator remains inaccurate, or objectives vary at a linear cumulative rate. Sublinear average residual requires the additional vanishing conditions stated in Corollary~\ref{cor:static}; it is not automatic in a continuously changing market. The optional residual error bound strengthens the interpretation to average distance from constrained stationary sets. Those sets can include nonglobal stationary points. Small windowed vector averages, by themselves, imply neither conclusion.

The bounded analysis domain is consequential. It can be checked directly in a finite controlled experiment, but its constants can grow with experiment length. An infinite-time conclusion requires a single uniformly bounded model family or an explicitly modified payoff/domain model. No unreported wealth clipping, stopping, or per-episode rescaling is silently used. Likewise, a fixed regularization defines a fixed target; sending the regularization to zero would require an additional approximation and uniformity analysis.

The mean recommendation is operationally separate from the stochastic policy optimized by the likelihood-gradient estimator. The observed context may be generated by mean execution without invalidating the conditional online telescope. Nevertheless, optimizing the stochastic-policy CPT objective does not prove optimality of the mean-execution policy or of terminal wealth.

\clearpage
\begin{thebibliography}{18}
\providecommand{\natexlab}[1]{#1}
\providecommand{\url}[1]{\texttt{#1}}
\expandafter\ifx\csname urlstyle\endcsname\relax
  \providecommand{\doi}[1]{doi: #1}\else
  \providecommand{\doi}{doi: \begingroup \urlstyle{rm}\Url}\fi

\bibitem[Arkes et~al.(2008)Arkes, Hirshleifer, Jiang, and
  Lim]{arkes2008reference}
Hal~R. Arkes, David Hirshleifer, Danling Jiang, and Sonya~S. Lim.
\newblock Reference point adaptation: Tests in the domain of security trading.
\newblock \emph{Organizational Behavior and Human Decision Processes},
  105\penalty0 (1):\penalty0 67--81, 2008.
\newblock \doi{10.1016/j.obhdp.2007.04.005}.

\bibitem[Arkes et~al.(2010)Arkes, Hirshleifer, Jiang, and Lim]{arkes2010cross}
Hal~R. Arkes, David Hirshleifer, Danling Jiang, and Sonya~S. Lim.
\newblock A cross-cultural study of reference point adaptation: Evidence from
  china, korea, and the {US}.
\newblock \emph{Organizational Behavior and Human Decision Processes},
  112\penalty0 (2):\penalty0 99--111, 2010.
\newblock \doi{10.1016/j.obhdp.2010.02.002}.

\bibitem[Aydore et~al.(2019)Aydore, Zhu, and Foster]{aydore2019dynamic}
Sergul Aydore, Tianhao Zhu, and Dean Foster.
\newblock Dynamic local regret for non-convex online forecasting.
\newblock In \emph{Advances in Neural Information Processing Systems},
  volume~32, 2019.

\bibitem[Fama and French(1993)]{fama1993common}
Eugene~F. Fama and Kenneth~R. French.
\newblock Common risk factors in the returns on stocks and bonds.
\newblock \emph{Journal of Financial Economics}, 33\penalty0 (1):\penalty0
  3--56, 1993.
\newblock \doi{10.1016/0304-405X(93)90023-5}.

\bibitem[Goda and Kitade(2023)]{goda2023unbiased}
Takashi Goda and Wataru Kitade.
\newblock Constructing unbiased gradient estimators with finite variance for
  conditional stochastic optimization.
\newblock \emph{Mathematics and Computers in Simulation}, 204:\penalty0
  743--763, 2023.
\newblock \doi{10.1016/j.matcom.2022.09.012}.
\newblock URL \url{https://arxiv.org/abs/2206.01991}.

\bibitem[Hallak et~al.(2021)Hallak, Mertikopoulos, and
  Cevher]{hallak2021regret}
Nadav Hallak, Panayotis Mertikopoulos, and Volkan Cevher.
\newblock Regret minimization in stochastic non-convex learning via a
  proximal-gradient approach.
\newblock In \emph{Proceedings of the 38th International Conference on Machine
  Learning}, pages 4008--4018. PMLR, 2021.

\bibitem[Hazan et~al.(2017)Hazan, Singh, and Zhang]{hazan2017efficient}
Elad Hazan, Karan Singh, and Cyril Zhang.
\newblock Efficient regret minimization in non-convex games.
\newblock In \emph{Proceedings of the 34th International Conference on Machine
  Learning}, pages 1433--1441. PMLR, 2017.

\bibitem[He and Yang(2019)]{he2019realization}
Xue~Dong He and Linan Yang.
\newblock Realization utility with adaptive reference points.
\newblock \emph{Mathematical Finance}, 29\penalty0 (2):\penalty0 409--447,
  2019.
\newblock \doi{10.1111/mafi.12182}.

\bibitem[Jie et~al.(2018)Jie, Prashanth, Fu, Marcus, and
  Szepesv{\'a}ri]{jie2018stochastic}
Cheng Jie, L.~A. Prashanth, Michael~C. Fu, Steven~I. Marcus, and Csaba
  Szepesv{\'a}ri.
\newblock Stochastic optimization in a cumulative prospect theory framework.
\newblock \emph{IEEE Transactions on Automatic Control}, 63\penalty0
  (9):\penalty0 2867--2882, 2018.
\newblock \doi{10.1109/TAC.2017.2777469}.

\bibitem[Kahneman and Tversky(1979)]{kahneman1979prospect}
Daniel Kahneman and Amos Tversky.
\newblock Prospect theory: An analysis of decision under risk.
\newblock \emph{Econometrica}, 47\penalty0 (2):\penalty0 263--291, 1979.
\newblock \doi{10.2307/1914185}.

\bibitem[Lepel and Barakat(2026)]{lepel2024policygradients}
Olivier Lepel and Anas Barakat.
\newblock Policy gradients for cumulative prospect theory in reinforcement
  learning.
\newblock \emph{Transactions on Machine Learning Research}, 2026.
\newblock URL \url{https://arxiv.org/abs/2410.02605v5}.
\newblock arXiv:2410.02605v5, camera-ready version, 6 September 2026.

\bibitem[Markowitz(1952)]{markowitz1952portfolio}
Harry Markowitz.
\newblock Portfolio selection.
\newblock \emph{The Journal of Finance}, 7\penalty0 (1):\penalty0 77--91, 1952.
\newblock \doi{10.2307/2975974}.

\bibitem[Prashanth et~al.(2016)Prashanth, Jie, Fu, Marcus, and
  Szepesv{\'a}ri]{prashanth2016cumulative}
L.~A. Prashanth, Cheng Jie, Michael~C. Fu, Steven~I. Marcus, and Csaba
  Szepesv{\'a}ri.
\newblock Cumulative prospect theory meets reinforcement learning: Prediction
  and control.
\newblock In \emph{Proceedings of the 33rd International Conference on Machine
  Learning}, pages 1406--1415. PMLR, 2016.

\bibitem[Qin et~al.(2022)Qin, Kong, and Yue]{qin2022realization}
Cong Qin, L.~H. Kong, and Xingye Yue.
\newblock Realization utility with path-dependent reference points.
\newblock \emph{SIAM Journal on Financial Mathematics}, 13\penalty0
  (3):\penalty0 912--948, 2022.
\newblock \doi{10.1137/21M1411457}.

\bibitem[Rhee and Glynn(2015)]{rhee2015unbiased}
Chang-han Rhee and Peter~W. Glynn.
\newblock Unbiased estimation with square root convergence for {SDE} models.
\newblock \emph{Operations Research}, 63\penalty0 (5):\penalty0 1026--1043,
  2015.
\newblock \doi{10.1287/opre.2015.1404}.

\bibitem[Tsybakov(2009)]{tsybakov2009introduction}
Alexandre~B. Tsybakov.
\newblock \emph{Introduction to Nonparametric Estimation}.
\newblock Springer, 2009.

\bibitem[Tversky and Kahneman(1992)]{tversky1992advances}
Amos Tversky and Daniel Kahneman.
\newblock Advances in prospect theory: Cumulative representation of
  uncertainty.
\newblock \emph{Journal of Risk and Uncertainty}, 5\penalty0 (4):\penalty0
  297--323, 1992.
\newblock \doi{10.1007/BF00122574}.

\bibitem[van~der Vaart(1998)]{vandervaart1998asymptotic}
Aad~W. van~der Vaart.
\newblock \emph{Asymptotic Statistics}.
\newblock Cambridge University Press, 1998.

\end{thebibliography}

\clearpage
\appendix

\section{Notation and Reproducibility Convention}\label{app:notation}
All conditional statistical results fix an episode-start context and parameter before drawing new simulation samples. Independence means conditional independence given that context. The online argument subsequently integrates those conditional statements. Limits in $M$ fix the episode and model; limits in $K$ require uniform constants across episodes. The minimax analysis, when included, is a separate fixed-parameter statistical experiment.

\begin{longtable}{@{}p{.24\textwidth}p{.70\textwidth}@{}}
\caption{Notation and information scope.}\label{tab:notation}\\
\toprule Symbol & Meaning \\\midrule\endfirsthead
\toprule Symbol & Meaning \\\midrule\endhead
$k,u,h$ & Episode index, within-episode date, and finite horizon.\\
$X_u,r_u,\omega_u$ & Wealth, reference in the same units, and simplex action.\\
$z_k,\cF_k$ & Nonreference initial context and complete pre-sampling information.\\
$P_k,\widetilde P_k$ & True market model and frozen fitted simulation model.\\
$J_k,\widetilde J_k$ & Their conditional stochastic-policy CPT objectives.\\
$Y_k$ & Trajectory-specific terminal wealth minus terminal reference.\\
$\Theta,q_{i,u},a_{i,u}$ & Parameter domain, features, and Dirichlet concentrations.\\
$a_0,\gamma$ & Normalization threshold and fixed parameter step size.\\
$G,H$ & Trajectory log-likelihood gradient and, where used, Hessian.\\
$M_2,M_H$ & Uniform score second moment and expected Hessian operator-norm bounds.\\
$C_1,C_2,C_3$ & Uniform derivative bounds for the fixed normalized weights.\\
$T_n,\mathcal T_\ell,\Delta_\ell$ & Plug-in statistic, its $n_\ell=N_02^\ell$ version, and antithetic difference.\\
$\mathsf L,p_\ell,a,c$ & Random precision level, its probabilities, tail exponent, and optional cap.\\
$Z,Z^{(c)},M_k$ & Untruncated/capped replication and count of independent complete replications.\\
$g_k,\widetilde g_k,\widehat g_k$ & True gradient, simulation gradient, and sampled update input.\\
$v_k,V_K$ & Adjacent-target value variation and its sum through $K-1$.\\
$d_k^J,\beta_k^{\rm sim}$ & Adjacent-target gradient variation and same-episode simulation bias.\\
$u_k,U_K,D_h$ & Nonreference value variation, its sum, and reference propagation modulus.\\
$Q_k,\Sta_k,\mathcal E_K$ & True normalized projected residual, its zero set, and average squared residual.\\
$\kappa,w,\rho$ & Optional residual error-bound constant and diagnostic window parameters.\\
$B$ & In the statistical rate analysis only, expected raw trajectory-call budget; not $M$.\\
\bottomrule
\end{longtable}

\paragraph{Proof interfaces.}
The reference proof needs the scalar recursion and a value modulus. The score proof needs a dominated differentiable trajectory law. The estimator proof needs Bernoulli empirical-survival moments, not continuity of the outcome law. The online proof accepts any direction estimator with the declared conditional MSE. The set-distance conclusions add the separate error bound. These interfaces distinguish primitive conditions from conclusions to be established.

\paragraph{Formal-verification scope.}
The source includes stable labels and \texttt{LEAN\_TARGET} comments. These identify statements for later formalization; they are not Lean proofs or evidence of an existing Lean build. Fundamental imported results would include conditional expectation and Fubini, differentiation under the integral, projection variational inequalities, iid laws of large numbers and CLT, and, in the minimax version, likelihood ratios and testing inequalities. The deterministic projection and telescoping lemmas can be formalized independently of CPT integration.

\section{Proofs for the Reference Mechanism}\label{app:reference-proofs}
\subsection{Proof of \texorpdfstring{\Cref{prop:path-dependence}}{the reference proposition}}
\begin{proof}
For path $A$, the intermediate reference is $(1-\eta_+)r_0+\eta_+H$. The assumed second-step branch gives
\[
 r_A=(1-\eta_-)((1-\eta_+)r_0+\eta_+H)+\eta_-C.
\]
For path $B$, the intermediate reference is $(1-\eta_-)r_0+\eta_-L$, so
\[
 r_B=(1-\eta_+)((1-\eta_-)r_0+\eta_-L)+\eta_+C.
\]
Subtracting cancels the $r_0$ terms and proves \eqref{eq:path-diff}. For a deterministic outcome $y$, endpoint normalization gives its CPT value as $v_+(y)-v_-(y)$: the relevant survival is one below the utility value and zero above it. This signed transformation is strictly increasing for the value functions specified here. Thus different deterministic relative outcomes have different CPT values.

Fix a wealth value $x$. The map $r\mapsto\Phi(r,x)$ is continuous at $r=x$ and has slopes $1-\eta_+$ for $r\le x$ and $1-\eta_-$ for $r>x$. If $r,r'$ lie on one side of $x$, its Lipschitz bound is immediate. If they straddle $x$, split the interval at $x$ and add the two bounds. Therefore
\[
 |\Phi(r,x)-\Phi(r',x)|\le(1-\eta_{\min})|r-r'|.
\]
Induction over the shared wealth path proves \eqref{eq:reference-contraction}. A common terminal wealth then gives $|Y-Y'|\le\chi_h|r-r'|$.

For $x>0$,
\[
 u'_{\alpha,\eps_v}(x)=\alpha x(x^2+\eps_v^2)^{\alpha/2-1}.
\]
Using $x\le\sqrt{x^2+\eps_v^2}$ and $\alpha-1\le0$ bounds this derivative by $\alpha\eps_v^{\alpha-1}$. Its limit at zero is zero, so composition with the positive/negative-part maps is continuously differentiable across zero and has the asserted global Lipschitz bound. Applying the two Lipschitz inequalities and adding them proves $\Omega(s)=(L_v^++L_v^-)s$.

Under Assumption~\ref{ass:reference-coupling}, use the same action and market draws for the two initial references. Write $V_\pm=v_\pm(Y)$, $V_\pm'=v_\pm(Y')$, and let $S_\pm,S_\pm'$ be their survivals. For any nonnegative $V,V'$ bounded by $B_v$, Fubini and the coupling give
\begin{equation}\label{eq:survival-coupling-inequality}
 \int_0^{B_v}|\Pp(V>z)-\Pp(V'>z)|\,dz
 \le \E\int_0^{B_v}|\1\{V>z\}-\1\{V'>z\}|\,dz
 =\E|V-V'|.
\end{equation}
The pointwise modulus bounds the sum of the two utility differences by $D_h(|r-r'|)$. Applying the mean value bound $|W(p)-W(q)|\le C_1|p-q|$ to the two CPT integrals proves \eqref{eq:reference-value-sensitivity}.

For the gradient claim, invoke the independently proved \Cref{thm:gradient-identity}. Centering the indicator in that identity only subtracts a deterministic scalar times a zero-mean score. Thus one can use
\[
 \varphi(\tau)=\int_0^{B_v}W_+'(S_+(z))\1\{V_+>z\}\,dz
 -\int_0^{B_v}W_-'(S_-(z))\1\{V_->z\}\,dz
\]
and $\nabla J=\E[G\varphi]$. The coupled wealth and action paths have the same features and the same $G$. For one sign,
\begin{align*}
 &|W'(S)\1\{V>z\}-W'(S')\1\{V'>z\}|\\
 &\hspace{1em}\le C_1|\1\{V>z\}-\1\{V'>z\}|+C_2|S-S'|.
\end{align*}
Multiply by $\norm{G}_2$, integrate, take expectations, and add the two signs. The first term is bounded by $C_1\E\norm{G}_2 D_h(|r-r'|)$ pathwise. The survival differences in the second term are deterministic under the fixed context; \eqref{eq:survival-coupling-inequality} bounds their integrated sum by $D_h(|r-r'|)$. Using $\E\norm{G}_2\le\sqrt{M_2}$ proves \eqref{eq:reference-gradient-sensitivity}. No bounded score or density for the outcome variable was used.
\end{proof}

\section{Proofs for the CPT Gradient and Estimator}\label{app:estimator-proofs}
\subsection{Proof of \texorpdfstring{\Cref{thm:gradient-identity}}{the CPT gradient identity}}
% LEAN_TARGET: dirichlet_moments_and_likelihood_domination
\begin{lemma}[Dirichlet likelihood moments and domination]\label{lem:dirichlet-moments}
Under the primitive compact-domain and parameter-independent-kernel conditions in Assumption~\ref{ass:bounded}, the trajectory likelihood is twice differentiable under integration against every bounded, parameter-independent trajectory functional. For every fixed real $p\ge1$ there is a finite uniform constant $M_p$ such that
\[
 \E_\theta\norm{G}_2^p\le M_p,\qquad
 \E_\theta\norm{H}_{\op}\le M_H<\infty,
 \qquad H=\nabla_\theta^2\log p_\theta(\tau).
\]
The constants may depend on $p,d,N,h,Q_f$ and the compact parameter neighborhood. No pointwise score bound is asserted.
\end{lemma}
\dependson{The exponential Dirichlet parametrization, bounded features on a bounded parameter neighborhood, a finite horizon, and a parameter-independent market kernel. Neither payoff differentiability in $Y$ nor the reference-isolated coupling is required.}
\begin{proof}
Choose a bounded open neighborhood $O$ of $\Theta$ whose closure is inside the allowed parameter domain, and let $B_\theta=\sup_{\theta\in\overline O}\norm\theta_2$. Uniformly in the state and $\theta\in O$,
\[
 \underline a=e^{-B_\theta Q_f}\le a_i\le e^{B_\theta Q_f}=\overline a,
 \qquad (N+1)\underline a\le A\le(N+1)\overline a.
\]
These inequalities follow from Cauchy--Schwarz in $\theta^\top q_i$. Set $t=\underline a/2$. Conditional on the pre-action state, the Dirichlet marginal negative moment is
\begin{equation}\label{eq:dirichlet-negative-moment}
 \E[\omega_i^{-t}\mid s]
 =\frac{\Gamma(A)\Gamma(a_i-t)}{\Gamma(a_i)\Gamma(A-t)}.
\end{equation}
The arguments of all gamma factors range over compact subsets of $(0,\infty)$, so this ratio has a finite uniform upper bound. The formula follows by integrating the Dirichlet density after replacing exponent $a_i-1$ by $a_i-t-1$ and taking the ratio of its normalizing constants. For $x\in(0,1]$, put $z=-\log x\ge0$. Since $z^p e^{-tz}$ is bounded on $[0,\infty)$, $|\log x|^p\le C_{p,t}x^{-t}$. Equation~\eqref{eq:dirichlet-negative-moment} therefore bounds every fixed conditional log moment.

The score contribution at one action is
\[
 G_u=\sum_i a_iq_i B_i,\qquad
 B_i=\log\omega_i-\psi_0(a_i)+\psi_0(A).
\]
All coefficients and digamma terms apart from the logarithms are uniformly bounded on the concentration domain. The inequality
$\norm{\sum_{j=1}^m x_j}_2^p\le m^{p-1}\sum_{j=1}^m\norm{x_j}_2^p$, together with conditional expectation and the finite sums over assets and dates, gives a uniform $M_p$. Dependence of later states on earlier actions does not require independent score contributions; conditional moment bounds suffice.

Let $\psi_1$ denote the trigamma function and $b=\sum_i a_iq_i$. Differentiation at a fixed trajectory, so each $q_i$ is fixed, gives
\begin{equation}\label{eq:dirichlet-hessian}
 H_u=\sum_i a_iq_iq_i^\top B_i
 +\psi_1(A)bb^\top
 -\sum_i\psi_1(a_i)a_i^2q_iq_i^\top.
\end{equation}
The bounded coefficients and conditional first log moments imply a uniform bound on $\E\norm{H_u}_{\op}$. Summing over $h$ dates gives $M_H$.

To justify rather than assume differentiation under the trajectory integral, define a reference probability law $\nu$ using the same market kernels and initial state, but independent-of-state Dirichlet action concentrations $b_*:=\underline a/2$ for every asset. Conditional on a complete trajectory, its likelihood ratio is a product over the finite horizon of factors
\[
 \frac{\Gamma(A)}{\prod_i\Gamma(a_i)}
 \frac{\Gamma(b_*)^{N+1}}{\Gamma((N+1)b_*)}
 \prod_i\omega_i^{a_i-b_*}.
\]
The normalizing factors and their first two parameter derivatives are uniformly bounded. Every exponent $a_i-b_*$ is at least $\underline a/2>0$. The first two derivatives of the product introduce only bounded coefficients and products of at most two logarithms. For each such power, $x^{\underline a/2}|\log x|^j$ is bounded on $(0,1]$. Thus the likelihood ratio and its first two derivatives admit finite constant envelopes with respect to the probability measure $\nu$. This argument also covers cross terms from different dates. Dominated differentiation now applies to every bounded trajectory functional. In particular
\begin{equation}\label{eq:second-likelihood-interface}
 \nabla_\theta^2\E_\theta F
 =\E_\theta[F(H+GG^\top)].
\end{equation}
The conditional market kernels cancel from the likelihood ratio because they have no direct parameter dependence; no density of the return law with respect to Lebesgue measure is needed. The same proof holds for each frozen simulator, with uniform constants from the common domain.
\end{proof}

\begin{proof}[Proof of the gradient identity]
Fix a true or simulation conditional trajectory law and a parameter. By the likelihood assumptions, for any bounded measurable trajectory functional $F$ without explicit parameter dependence,
\begin{equation}\label{eq:likelihood-interface}
 \nabla_\theta\E_\theta F=\E_\theta[FG].
\end{equation}
In particular $\E G=\nabla 1=0$ and, for each threshold $z$,
\[
 \nabla S_\pm(z)=\E[\1\{v_\pm(Y)>z\}G]
 =\E[(\1\{v_\pm(Y)>z\}-S_\pm(z))G].
\]
Since $\norm{\nabla S_\pm(z)}_2\le\sqrt{M_2}$ and $|W_\pm'|\le C_1$, the derivative integrands have an integrable bound over the finite interval $[0,B_v]$. The chain rule and differentiation under the integral give \eqref{eq:gradient-identity}. Fubini is valid because $|\psi|\le2B_vC_1$ and $\E\norm{G}_2<\infty$. This also yields \eqref{eq:gradient-bound}. The path recursion is held fixed as a measurable trajectory functional; no derivative of the historical starting state is taken.

For completeness, twice dominated likelihood differentiation in Lemma~\ref{lem:dirichlet-moments} gives
\[
 \nabla^2 S_\pm(z)=\E[\1\{v_\pm(Y)>z\}(H+GG^\top)].
\]
Therefore $\norm{\nabla^2S_\pm(z)}_{\op}\le M_H+M_2$, and the first-derivative bound above is $\norm{\nabla S_\pm(z)}_2\le\sqrt{M_2}$. The Hessian chain rule for either weight yields
\[
 \nabla^2 W_\pm(S_\pm)
 =W_\pm''(S_\pm)\nabla S_\pm\nabla S_\pm^\top
 +W_\pm'(S_\pm)\nabla^2S_\pm.
\]
Its operator norm is at most $C_2M_2+C_1(M_2+M_H)$ at every threshold. Integrating over the two intervals of length $B_v$ gives the bound \eqref{eq:derived-smoothness}. For $\theta,\theta'\in\Theta$, the line segment between them lies in $\Theta$; integration of this Hessian along the segment proves
$\norm{\nabla J(\theta')-\nabla J(\theta)}_2\le L\norm{\theta'-\theta}_2$.
The reference recursion and even atoms in $Y$ cause no difficulty because each threshold indicator is a bounded parameter-independent trajectory functional.
\end{proof}

\subsection{Proof of \texorpdfstring{\Cref{prop:debias}}{the antithetic estimator proposition}}
\begin{proof}
Fix the episode-start context, parameter, and simulation law. Write $S$ for its true survival and condition initially on the outer trajectory. For $I\in\{0,1\}$,
\[
 f'(p,I)=W''(p)(I-p)-W'(p),\qquad
 f''(p,I)=W'''(p)(I-p)-2W''(p).
\]
Hence $|f''|\le H_f$. At every fixed threshold, $\widehat S_n$ is an average of independent Bernoulli variables of mean $S$, independently of the outer trajectory. Consequently
\[
 \E(\widehat S_n-S)=0,\qquad
 \E(\widehat S_n-S)^2=S(1-S)/n\le1/(4n).
\]
Taylor's theorem with integral remainder implies
\[
 |\E[f(\widehat S_n,I)\mid\tau^{(0)}]-f(S,I)|
 \le H_f/(8n).
\]
Integrating over the two intervals of length $B_v$, multiplying by the outer score, and applying $\E\norm{G_0}_2\le\sqrt{M_2}$ proves
$\norm{\E T_n-\E T_\infty}_2\le A_b/n$.
Theorem~\ref{thm:gradient-identity} identifies $\E T_\infty=\widetilde g_k$. Also $\norm{\mathcal T_0}_2\le2B_vC_1\norm{G_0}_2$, giving the base-level bound $C_0$.

For the antithetic bound, set $n=n_\ell=2m$. At a fixed threshold, let $x$ and $y$ be the two half-sample survival means. Their full mean is exactly $(x+y)/2$. Taylor expansion about that midpoint gives
\begin{equation}\label{eq:midpoint-remainder}
 \left|f\left(\frac{x+y}2,I\right)-\frac{f(x,I)+f(y,I)}2\right|
 \le\frac{H_f}{8}(x-y)^2.
\end{equation}
To bound the fourth moment, pair Bernoulli variables from the two halves and let $D_j=B_j-C_j$. These are independent, centered, and satisfy
$\E D_j^2=\E D_j^4=2S(1-S)\le1/2$. Expanding the fourth power and eliminating odd cross terms yields
\begin{align}
 \E|x-y|^4
 &=\frac{m\E D_1^4+3m(m-1)(\E D_1^2)^2}{m^4}\notag\\
 &\le\frac{3}{4m^2}=\frac3{n^2}.
 \label{eq:bernoulli-fourth-moment}
\end{align}
Let $R_+(z),R_-(z)$ denote the scalar midpoint remainders. Cauchy--Schwarz over the disjoint union of two intervals gives
\[
 \left|\int R_+-\int R_-\right|^2
 \le2B_v\left(\int|R_+|^2+\int|R_-|^2\right).
\]
The inner-sample bound is uniform in the outer indicators. Condition on the outer trajectory and use \eqref{eq:midpoint-remainder}--\eqref{eq:bernoulli-fourth-moment}; then take the outer expectation. This gives
\[
 \E\norm{\Delta_\ell}_2^2
 \le M_2(2B_v)(2B_v)\frac{H_f^2}{64}\frac3{n_\ell^2}
 =D2^{-2\ell}.
\]
No independence across thresholds or between gain/loss indicators is required.

The marginal distribution of each half-statistic is that of $\mathcal T_{\ell-1}$, despite sharing the outer trajectory. Therefore
$\E\Delta_\ell=\E \mathcal T_\ell-\E \mathcal T_{\ell-1}$ for $\ell\ge1$. The moment bounds imply
\[
 \sum_{\ell=0}^\infty\E\norm{\Delta_\ell}_2
 \le\sqrt{C_0}+\sqrt D\sum_{\ell=1}^\infty2^{-\ell}<\infty.
\]
Absolute integrability justifies summing expectations for the independent randomized level:
\[
 \E Z=\sum_{\ell=0}^\infty\E\Delta_\ell
 =\lim_{\ell\to\infty}\E \mathcal T_\ell=\widetilde g_k.
\]
Similarly, nonnegative summation gives
\[
 \E\norm{Z}_2^2
 =\sum_{\ell=0}^\infty\frac{\E\norm{\Delta_\ell}_2^2}{p_\ell}
 \le\frac{C_0+D\sum_{\ell=1}^\infty2^{-(2-a)\ell}}{1-2^{-a}}
 =C_{\db}.
\]
Because one replication uses exactly $1+N_02^{\mathsf L}$ raw trajectories, summing its cost gives $C_{\rm cost}$ in \eqref{eq:uniform-db-constant}. The two convergent geometric series respectively require $a<2$ and $a>1$. Multiplication by a level factor $O(1+\ell)$ for sorting leaves the cost series convergent.

For a finite cap let $P_c=\sum_{j=0}^cp_j$. Its distribution is $q_\ell=p_\ell/P_c$ on $0,\ldots,c$. Then
\[
 \E Z^{(c)}=\sum_{\ell=0}^c\E\Delta_\ell=\E \mathcal T_c,
 \qquad
 \E\norm{Z^{(c)}}_2^2=P_c\sum_{\ell=0}^c
 \frac{\E\norm{\Delta_\ell}_2^2}{p_\ell}\le C_{\db}.
\]
The weak-bias bound with $n_c=N_02^c$ proves the remaining claim. This uses the specified normalized capped distribution, not an unspecified operation that clips a sampled level while retaining its old probability.
\end{proof}

\subsection{Proof of \texorpdfstring{\Cref{thm:variance}}{the variance theorem}}
\begin{proof}
Fix the conditional context. For independent complete replications, cross-covariances of centered outputs vanish. Thus
\[
 \tr\Var\left(M^{-1}\sum_{b=1}^MZ_b\right)
 =M^{-1}\tr\Var(Z_1)\le C_{\db}/M.
\]
The mean is $\widetilde g_k$ without truncation and $\E \mathcal T_c$ with cap $c$. Decompose
\[
 \E\norm{\widehat g-g_k}_2^2
 =\tr\Var(\widehat g)+\norm{\E\widehat g-g_k}_2^2.
\]
The deterministic conditional bias is a sum of simulation bias and truncation bias. Applying $\norm{a+b}_2^2\le2\norm a_2^2+2\norm b_2^2$ proves \eqref{eq:total-mse}.

For the plug-in estimator, denote its whole inner sample by $\mathcal I$. Conditional on $\mathcal I$, the $m$ outer terms are independent and each has second moment at most $C_0$. Let $m(\mathcal I)$ be the conditional mean of one term. The trace form of total variance yields
\begin{equation}\label{eq:plugin-total-variance}
 \tr\Var(\widehat g_{n,m}^{\pl})
 \le C_0/m+\E\norm{m(\mathcal I)-\E m(\mathcal I)}_2^2.
\end{equation}
Replacing one inner trajectory changes either empirical survival by at most $1/n$ at every threshold. Since $|f'(p,I)|\le C_1+C_2$, the norm of the resulting conditional-mean change is at most
$2B_v(C_1+C_2)\sqrt{M_2}/n$.
The vector Efron--Stein inequality follows by applying its scalar version to each coordinate and summing. With $\mathcal I^{(i)}$ formed by independent replacement of inner sample $i$, it gives
\[
 \E\norm{m(\mathcal I)-\E m(\mathcal I)}_2^2
 \le\frac12\sum_{i=1}^n\E\norm{m(\mathcal I)-m(\mathcal I^{(i)})}_2^2
 \le\frac{2B_v^2(C_1+C_2)^2M_2}{n}.
\]
The mean bias relative to $\widetilde g_k$ is at most $A_b/n$ by \Cref{prop:debias}. Add simulation bias, square with the same two-term inequality, and combine with \eqref{eq:plugin-total-variance}. This proves \eqref{eq:plugin-mse} while retaining the shared-inner-sample dependence.
\end{proof}

\subsection{Proof of \texorpdfstring{\Cref{thm:clt}}{the fixed-episode CLT}}
\begin{proof}
For almost every fixed conditional context, the centered vectors $Z_b-\widetilde g_k$ are iid with finite covariance. For each deterministic $t\in\R^d$, the scalar variables $t^\top(Z_b-\widetilde g_k)$ are iid with zero mean and finite variance $t^\top\Sigma_kt$. The scalar iid CLT and Cramer--Wold theorem yield \eqref{eq:clt}, including singular covariance matrices. These probability limit theorems are applied to a fixed conditional distribution, not to a changing sequence of models.

Each covariance entry has finite first moment because
$\E|Z_iZ_j|\le(\E Z_i^2\E Z_j^2)^{1/2}<\infty$. The iid strong law applied to entries and to sample means proves consistency of the empirical covariance. If $\Sigma_k$ is positive definite, the inverse square-root map is continuous in a neighborhood of $\Sigma_k$. With probability tending to one the empirical covariance lies in that neighborhood; defining the standardized statistic arbitrarily on its complement does not affect the limit. Slutsky's theorem proves the studentized claim. The proof needs no $(2+\xi)$ moment. If the simulator is inaccurate, centering at $g_k$ instead adds the nonvanishing deterministic shift $\sqrt M(\widetilde g_k-g_k)$ and is not authorized.
\end{proof}

\section{Proofs for Stationary-Set Drift and Dynamic Regret}\label{app:regret-proofs}
\subsection{Proof of \texorpdfstring{\Cref{lem:set-drift}}{the stationary-set drift lemma}}
\begin{proof}
The map $d(g)$ is the Euclidean projection of $g/a_0$ onto the closed unit ball. Nonexpansiveness gives
\begin{equation}\label{eq:normalization-lipschitz}
 \norm{d(g)-d(g')}_2\le a_0^{-1}\norm{g-g'}_2.
\end{equation}
Nonexpansiveness of $\Pi_\Theta$ and the common step size therefore imply
\[
 \sup_{\theta\in\Theta}\norm{Q_{k+1}(\theta)-Q_k(\theta)}_2
 \le d_k^J/a_0.
\]
For $\theta\in\Sta_k$, $Q_k(\theta)=0$. The new-episode error bound gives
$\dist(\theta,\Sta_{k+1})\le\kappa d_k^J/a_0$.
Take the supremum over the old set. Interchanging the two episodes gives the reverse directed Hausdorff inequality with the same uniform $\kappa$. Both sets are nonempty compact subsets of $\Theta$, so their Hausdorff distance is well-defined and bounded as claimed.
\end{proof}

\subsection{Projected ascent inequality}
% LEAN_TARGET: normalized_projected_one_step
\begin{lemma}[One-step normalized projected ascent]\label{lem:projected-ascent}
Let $J$ be $L$-smooth on nonempty closed convex $\Theta$, $x\in\Theta$, $g=\nabla J(x)$, and $\norm g_2\le B_g$. For any finite vector $\widehat g$, define
\[
 x^+=\Pi_\Theta(x+\gamma d(\widehat g)),\qquad
 Q(x)=\gamma^{-1}\{\Pi_\Theta(x+\gamma d(g))-x\}.
\]
If $a_0,\gamma>0$ and $\gamma L\le a_0$, with $C_g=\max(B_g,a_0)$,
\begin{equation}\label{eq:projected-ascent}
 J(x^+)-J(x)
 \ge\frac{a_0\gamma}{8}\norm{Q(x)}_2^2
 -\frac{5C_g\gamma}{4a_0^2}\norm{\widehat g-g}_2^2.
\end{equation}
\end{lemma}
\dependson{Only the deterministic smoothness and projection conditions in the statement. No unbiasedness, bounded sample score, alignment condition, or stationary-set error bound is needed.}
\begin{proof}
Put $s=x^+-x$, $c=\max(\norm g_2,a_0)$, and $\delta=d(\widehat g)-d(g)$. The projection variational inequality, tested at the feasible point $x$, gives
\[
 \ip{d(\widehat g)}{s}\ge\norm s_2^2/\gamma.
\]
Since $g=c\,d(g)$, the smoothness lower bound yields
\begin{align}
 J(x+s)-J(x)
 &\ge\ip g s-\tfrac L2\norm s_2^2\notag\\
 &\ge(c/\gamma-L/2)\norm s_2^2-c\norm\delta_2\norm s_2\notag\\
 &\ge(3c/(4\gamma)-L/2)\norm s_2^2-c\gamma\norm\delta_2^2\notag\\
 &\ge\frac{c}{4\gamma}\norm s_2^2-c\gamma\norm\delta_2^2.
 \label{eq:actual-displacement-bound}
\end{align}
The third line uses $cuv\le c u^2/(4\gamma)+c\gamma v^2$ with $u=\norm s_2$, $v=\norm\delta_2$; the fourth uses $\gamma L\le a_0\le c$.

Let $t=\Pi_\Theta(x+\gamma d(g))-x=\gamma Q(x)$. Nonexpansiveness gives $\norm{s-t}_2\le\gamma\norm\delta_2$. Consequently
\[
 \norm s_2^2\ge\tfrac12\norm t_2^2-\gamma^2\norm\delta_2^2.
\]
Substitute in \eqref{eq:actual-displacement-bound} to obtain
\[
 J(x^+)-J(x)\ge\frac{c\gamma}{8}\norm{Q(x)}_2^2
 -\frac{5c\gamma}{4}\norm\delta_2^2.
\]
Now $a_0\le c\le C_g$ and \eqref{eq:normalization-lipschitz} imply \eqref{eq:projected-ascent}. In particular the inequality remains valid when projection removes the whole feasible step.
\end{proof}

\subsection{Proof of \texorpdfstring{\Cref{thm:dynamic-regret}}{the online residual theorem}}
\begin{proof}
Apply \Cref{lem:projected-ascent} episode by episode, condition on $\cF_k$, and then take total expectation. With $\Delta J_k=J_k(\theta_{k+1})-J_k(\theta_k)$,
\begin{equation}\label{eq:one-step}
 \frac{a_0\gamma}{8}\E\norm{Q_k(\theta_k)}_2^2
 \le\E\Delta J_k+\frac{5C_g\gamma}{4a_0^2}\E e_k^2.
\end{equation}
For every realized sequence of objectives and iterates,
\begin{align}
 \sum_{k=1}^K\Delta J_k
 ={}&J_K(\theta_{K+1})-J_1(\theta_1)\notag\\
 &+\sum_{k=1}^{K-1}
 [J_k(\theta_{k+1})-J_{k+1}(\theta_{k+1})]
 \le2B_J+V_K.
 \label{eq:pathwise-telescope}
\end{align}
The identity is purely algebraic; endogenous generation of the next context does not invalidate it. Summing \eqref{eq:one-step}, using \eqref{eq:pathwise-telescope}, and dividing by $a_0\gamma K/8$ proves \eqref{eq:main-bound}. The constants $8$ and $10$ follow from that division. \Cref{thm:variance} supplies \eqref{eq:algorithm-mse-substitution} for the implemented estimator.

For the window statement, put $\alpha_{k,j}=\rho^j/A_k$. Jensen's inequality gives
\[
 \norm{\overline Q_k}_2^2
 \le\sum_{j=0}^{\min(w-1,k-1)}\alpha_{k,j}
 \norm{Q_{k-j}(\theta_{k-j})}_2^2.
\]
Each residual occurs in at most $w$ windows and each coefficient is at most one. Summing gives $\Reg_{\rho,w}^Q(K)\le w\sum_k\norm{Q_k(\theta_k)}_2^2$. A growing window must retain this factor.

Finally split the adjacent-objective difference at the intermediate target $J(P_{k+1},z_{k+1},r_k;\theta)$. Triangle inequality and \eqref{eq:reference-value-sensitivity} give
\[
 v_k\le u_k+C_1D_h(|r_{k+1}-r_k|).
\]
Sum this pathwise inequality to obtain \eqref{eq:reference-budget}. It retains market changes, starting wealth, previous positions, and market initial state in $u_k$; none is silently dropped. No bounded-Lipschitz distance is presumed to control an unbounded score kernel.
\end{proof}

\subsection{Proof of \texorpdfstring{\Cref{cor:static}}{the static-limit corollary}}
\begin{proof}
For $J_k\equiv J$, $V_K=0$. The first term of \eqref{eq:main-bound} tends to zero for fixed positive $a_0,\gamma$, and the second tends to zero by the average-MSE condition. For nonstatic objectives the additional condition $\E V_K=o(K)$ has the same consequence. In exact untruncated simulation, $e_k^2=C_{\db}/M_k$ is admissible. If $M_k\to\infty$ almost surely, then $1/M_k\to0$ almost surely and $0\le1/M_k\le1$. Dominated convergence gives $\E[1/M_k]\to0$, and its Cesaro mean tends to zero. These statements concern averages of expected squared residuals; they do not imply pointwise or almost-sure convergence of the iterates.
\end{proof}

\subsection{Proof of \texorpdfstring{\Cref{cor:set-tracking}}{the set-tracking corollary}}
\begin{proof}
For every analyzed iterate the optional error bound gives
$\dist^2(\theta_k,\Sta_k)\le\kappa^2\norm{Q_k(\theta_k)}_2^2$.
Take expectations and average. Substitution of \eqref{eq:main-bound} proves the finite-time assertion; the limiting assertion follows from \Cref{cor:static}. The proof does not use a converse to Jensen or the stationary-set drift lemma.
\end{proof}

\section{Proof of the Minimax Lower Bound}\label{app:minimax-proof}
\subsection{Proof of \texorpdfstring{\Cref{thm:minimax}}{the minimax theorem}}
\begin{proof}
We specify an actual portfolio-trajectory experiment before applying any statistical testing inequality.

\paragraph{Step 1: a realizable submodel.}
Use one stock and cash, $h=1$, $c_{\rm tc}=0$, $X_0=r_0=1$, and fixed $\eta_+\in(0,1)$; fix any $\eta_-\in[0,1]$. Set $q_{\rm risky}=e_1$ and $q_{\rm cash}=0$, where $e_1$ is the first coordinate vector in $\R^d$. Thus the risky weight is $U\sim\operatorname{Beta}(a,1)$ with $a=e^{\theta_1}$, and the remaining parameters do not affect the law. Let the stock return be $R=r_*H$, where $H\sim\operatorname{Bernoulli}(p)$ is independent of $U$ and $r_*>0$ is fixed and small enough for the declared bounds. A complete oracle trajectory records the action and return, so it reveals $H$.

The reference update is always on the gain or equality branch:
\[
 X_1=1+r_*UH,\qquad
 r_1=1+\eta_+r_*UH,\qquad
 Y=cUH,\quad c=(1-\eta_+)r_*>0.
\]
This is exactly the wealth/reference recursion in the manuscript. States and values are bounded, the feature map is bounded, and the compact parameter domain supplies the uniform likelihood bounds from \Cref{lem:dirichlet-moments}. A fixed compact interval $p\in I\subset(0,1)$ changes none of these bounds. Choose the common class constants to dominate this submodel's constants, as required in Assumption~\ref{ass:oracle}. The loss utility vanishes. We have not postulated a distribution of abstract gradient vectors.

\paragraph{Step 2: separation of genuine CPT gradients.}
Let $v=v_+$ be the fixed smooth gain utility and put $m(t)=c\,v'(ct)$ for $t\in[0,1]$. Here $m(t)>0$ for $t>0$ and $m$ is bounded. The substitution $z=v(ct)$ in the CPT survival integral gives
\begin{equation}\label{eq:submodel-cpt}
 J_p(\theta)=\int_0^1 m(t)W_+\{p(1-t^{e^{\theta_1}})\}\,dt.
\end{equation}
Indeed, for $0<t<1$, $\Pp\{v(Y)>v(ct)\}=p\Pp(U>t)=p(1-t^{e^{\theta_1}})$. Above $v(c)$ the survival is zero and $W_+(0)=0$. Differentiating the bounded-integrable expression at $\theta_0=0$ yields
\begin{equation}\label{eq:submodel-gradient}
 \nabla J_p(0)=g(p)e_1,\qquad
 g(p)=p\int_0^1m(t)W_+'(p(1-t))(-t\log t)\,dt.
\end{equation}
The functions $t^a|\log t|$ are integrable with uniform bounds on a compact positive range of $a$, so the differentiation is legitimate. Put
\[
 A_v=\int_0^1m(t)(-t\log t)\,dt>0,
 \qquad w_0=W_+'(0)>0.
\]
Differentiate once in $p$, using $|W_+''|\le C_2$:
\begin{align*}
 g'(p)&=\int_0^1m(t)(-t\log t)
 [W_+'(p(1-t))+p(1-t)W_+''(p(1-t))]dt\\
 &\ge (w_0-2C_2p)A_v.
\end{align*}
Choose $p_{\max}\in(0,1/4]$ with $2C_2p_{\max}\le w_0/2$, possible also when $C_2=0$. Define $p_*=p_{\max}/2$, $I=[p_*/2,3p_*/2]$, and $m_*=w_0 A_v/2>0$. Then $g'(p)\ge m_*$ throughout $I$. This nonzero derivative proves local identifiability of the gradient functional, rather than inferring it just from identifiability of $p$.

\paragraph{Step 3: information in a stopped adaptive experiment.}
For $B\ge1$ choose $p_\pm=p_*\pm\delta$, $\delta=a_*/\sqrt B$, where the fixed $a_*>0$ will be sufficiently small. Taking $a_*\le p_*/2$ keeps $p_\pm\in I$. A one-call Bernoulli relative entropy obeys
\[
 \operatorname{kl}(p_+\Vert p_-)
 \le\frac{(p_+-p_-)^2}{p_-(1-p_-)}
 \le C_I\delta^2,
 \qquad C_I=\frac{4}{(p_*/2)(1-3p_*/2)}.
\]
The first inequality follows from $\log x\le x-1$ applied to each of the two Bernoulli likelihood terms. At any policy query the conditional action distribution is the same under $p_+$ and $p_-$; the only differing random kernel is the Bernoulli return. Thus the relative entropy per fresh trajectory is exactly this Bernoulli relative entropy even when the query depends on previous observations.

For completeness, include all independent algorithmic randomness in the transcript and let $N=N_{\rm call}$. Conditional query and stopping rules are the same functions of the observed history under the two models. On an observed stopped history of length $n$, the likelihood ratio is therefore
\[
 \prod_{j=1}^n
 \left(\frac{p_+}{p_-}\right)^{H_j}
 \left(\frac{1-p_+}{1-p_-}\right)^{1-H_j}.
\]
This follows first on each event $\{N=n\}$ by multiplying conditional kernels; the common action, randomization, and stopping factors cancel. Both models have finite expected $N$, hence $N<\infty$ almost surely and these events exhaust the stopped experiment. The likelihood increments are bounded because $p_\pm$ lie in the fixed interior interval $I$. Their absolute stopped sum has expectation at most a finite constant times $\E_+N$. We may consequently interchange its series and expectation. Since $\{N\ge j\}$ is measurable before the $j$th fresh return,
\begin{align}
 \operatorname{KL}(\mathbb P_+^{\rm tr}\Vert\mathbb P_-^{\rm tr})
 &=\sum_{j\ge1}\E_+\left[\1\{N\ge j\}
 \log\frac{p_+^{H_j}(1-p_+)^{1-H_j}}
 {p_-^{H_j}(1-p_-)^{1-H_j}}\right]\notag\\
 &=\E_+N\,\operatorname{kl}(p_+\Vert p_-)
 \le BC_I\delta^2=C_Ia_*^2.\label{eq:stopped-kl}
\end{align}
This proves the information bound for an expected budget and adaptive stopping, not only for a deterministic sample size. Choose $a_*$ still smaller so that $C_Ia_*^2\le1/8$. Pinsker's inequality gives total variation at most $1/4$.

\paragraph{Step 4: estimation-to-testing reduction.}
By the mean value theorem, the Euclidean target separation is
\[
 D_B=\norm{g(p_+)e_1-g(p_-)e_1}_2
 \ge 2m_*\delta.
\]
For any estimator form the test selecting the nearer of these two targets. On a testing error its distance from the true target is at least $D_B/2$ by the triangle inequality. The sum of the two test error probabilities is at least one minus total variation: for a decision event $A$, this sum is $\mathbb P_+(A^c)+\mathbb P_-(A)=1-[\mathbb P_+(A)-\mathbb P_-(A)]$. It follows that
\begin{align*}
 \max_{s\in\{+,-\}}\E_s\norm{\widehat g-g(p_s)e_1}_2^2
 &\ge\frac{D_B^2}{8}
 [1-\operatorname{TV}(\mathbb P_+^{\rm tr},\mathbb P_-^{\rm tr})]\\
 &\ge\frac{3m_*^2a_*^2}{8B}.
\end{align*}
Every procedure obeying the class-wide expected budget also obeys it on this two-point submodel. Taking the infimum over such procedures proves the class lower bound with $c_{\rm lb}=3m_*^2a_*^2/8>0$. The testing inequalities are classical \citep{tsybakov2009introduction}; the model embedding and stopped information calculation above provide their required inputs here.

\paragraph{Step 5: matching cost-calibrated upper bound.}
For every $P\in\mathfrak C$, the common primitive bounds give uniform $C_{\db}$ and the same $C_{\rm cost}$ in \Cref{prop:debias}. Run $M=\lfloor B/C_{\rm cost}\rfloor$ complete independent, untruncated replications at $\theta_0$. Their total number of trajectories is a permissible stopping time when the level and required number of remaining draws are part of the algorithmic state. Its expected cost is $MC_{\rm cost}\le B$. Exact sampling and \Cref{thm:variance} give
\[
 \sup_{P\in\mathfrak C}\E_P\norm{\widehat g-g_P}_2^2
 \le C_{\db}/M\le 2C_{\db}C_{\rm cost}/B
\]
for $B\ge2C_{\rm cost}$. This proves both sides on the same model class, with the same target, loss, oracle, and expected raw-trajectory budget.
\end{proof}


\section{Experimental Reproducibility Details}\label{app:experiments}
The accompanying \texttt{verification\_experiments.py} executes the statistics and recursions in the manuscript. It requires Python and NumPy; no external price feed, language model, or user-behavior dataset is accessed. The seed, all parameters, raw trajectory counts, estimator summaries, and episode records are saved to JSON and CSV. Table entries are produced from those saved outputs, not from prescribed performance curves. The estimator/learning seed is 20260911; the real execution stream uses seed $88177$. Level diagnostics use $1600$ independent replications per level, and batch diagnostics use $8192$ capped replications. Every batch reuses the same saved pool only across the reported group sizes; the replay uses fresh learning draws. No untruncated simulation is represented by a capped numerical run.


\paragraph{Analytic bias in the polynomial-weight diagnostic.}
Write $b=1-2\eps_p$, $D_w=(1-\eps_p)^3-\eps_p^3$. Then
\[
 A_0=3\eps_p^2b/D_w,\qquad A_1=6\eps_pb^2/D_w,
 \qquad A_2=3b^3/D_w.
\]
At $\theta=0$ in the pure-gain diagnostic, $S(v_+(ct))=1-t$ and
$\E[(1+\log U)\1\{U>t\}]=-t\log t$, by direct integration of $1+\log u$ from $t$ to $1$. Outer/inner independence and $\E G=0$ remove the term containing only $W'(\widehat S)\widehat S$ from $\E T_n$. Since a Bernoulli empirical mean has $\E\widehat S=S$ and $\E\widehat S^2=S^2+S(1-S)/n$, one has
\[
 \E W'(\widehat S)=W'(S)+A_2S(1-S)/n.
\]
Multiplying by $\E[G\1\{U>t\}]$, substituting $z=v_+(ct)$, and integrating proves \eqref{eq:diagnostic-bias}. All integrals are finite; for the power value at zero the factor is bounded in integral by a constant times $t^{\alpha_+}|\log t|$, and for the smooth value it is continuous apart from the harmless logarithmic endpoint factor. A cap $c$ targets $T_{N_02^c}$ in sample-count notation, giving the reported residual bias. The formulas provide an independent population comparison for the actual sampled statistics, not substitute data for them.

\paragraph{Exact integration of finite-sample statistics.}
For one sign, form the sorted unique union of zero, the inner utility observations, and the outer utility. On each adjacent interval all empirical survival and outer-indicator values are constant. Multiply the interval length by $W'(\widehat S)(I_0-\widehat S)$ and sum. Above the largest breakpoint both quantities are zero, so the integrand is zero because $f(0,0)=0$. Repeat for the loss side and multiply the difference by the shared score. The two half-statistics use the same outer value and score but their own inner breakpoints.

\paragraph{Independent numerical gradient check.}
The one-period replay has one risky asset with return $+r_*$ or $-r_*$ and $U\sim\operatorname{Beta}(e^\theta,1)$. Fixing the context $(X,r)$, a positive relative-outcome threshold $y$ is equivalent to terminal wealth exceeding $r+y/(1-\eta_+)$; a loss threshold is equivalent to terminal wealth being below $r-y/(1-\eta_-)$. For a return $R\ne0$, the corresponding action threshold is $(\text{wealth threshold}/X-1)/R$, clipped to $[0,1]$. Its CDF is $t^{e^\theta}$ and its parameter derivative in the interior is $e^\theta t^{e^\theta}\log t$. Combining the two return branches gives $S_\pm$ and $\nabla S_\pm$ directly. We integrate $W(S_\pm)$ and $W'(S_\pm)\nabla S_\pm$ over utility thresholds, splitting at endpoint breakpoints. Central differences of the independently integrated objective cross-check the gradient. Doubling the quadrature order supplies the reported numerical stability diagnostic; it is not a rigorous integration-error certificate.

\paragraph{Replay step-size domain.}
For the executed controlled replay, $h=1$, $c_{\rm tc}=0$, $|R|\le0.04$, $K=36$, and $\Theta=[-2,2]$. Mean or sampled actions lie in $[0,1]$, so all replay starting and one-step hypothetical wealth values are bounded by $X_{\max}=1.04^{K+1}$. References are convex combinations of earlier wealth and the initial reference, so the same upper bound applies. Nonnegativity gives $|Y|\le X_{\max}$. Set $B_v=\lambda v_+(X_{\max})$ since both exponents and smoothing scales coincide in this verification. The scalar Beta score is $G=1+e^\theta\log U$; $-e^\theta\log U$ is unit exponential, hence $\E G^2=1$ and $\E|\partial_\theta G|=1$. Thus $L=2B_v(C_2+2C_1)$ is a valid conservative bound from the likelihood Hessian calculation, and the implemented $a_0=1$, $\gamma=1/(2L)$ satisfies $\gamma L\le a_0$.

\paragraph{Scope of formal and numerical verification.}
The manuscript supplies mathematical proofs under explicit hypotheses. Compilation checks syntax, labels, and document consistency. The numerical checks test specified computations on finite instances. Neither substitutes for a Lean proof of integration, conditional independence, asymptotic limits, or the minimax experiment. No existing local Lean algebra checks are presented as a formalization of the new manuscript.
\end{document}
